\clearpage
\part*{Part IV. Spacetime, Symmetry And The Dynamics Of The Read}
\label{atlaspart:04-spacetime-and-symmetry}
\addcontentsline{toc}{part}{Part IV. Spacetime, Symmetry And The Dynamics Of The Read}
\section*{Introduction to Part IV}\label{narr:part:III}
\addcontentsline{toc}{section}{Introduction to Part IV}
Nothing in the axiom mentions space. Part IV asks where the order of events, the rotations, the Lorentz group, the boosts, spin and the discrete reflections come from in a construction whose only material is finitely many entangled pairs. Its answer rests on the table of nouns of chapter 1-7 more than on any other single result, because the symmetry chapters read the structure of the construction at the words the table fixes: the derived directions, the Hermitian split, the open-end label and the order parameter. The Part is short, eight chapters, and most of its results are conditional theorems whose conditions are listed as the first sentence inside the box.

The order of the chapters follows the order in which the objects of the literature are usually met, but the reason for it here is one of dependence. The order coordinate is read first because it is the object every other chapter of the Part refers to, the count against which the group moves and the transfers are measured. The rotations come second because they are the symmetries of a finite set of directions that chapter 1-7 already derives, and the Lorentz group third because it is built from the rotations by adjoining the boosts. The boost chapter and the spin chapter stand next as the two places where the construction stops short, each an open result with its parting named. The discrete reflections, the theorem on the order parameter and the empty chapter on causal structure close the Part. A reader who wants only what is proved can read chapters 3-1, 3-2, 3-3, 3-6 and 3-7 as a chain and treat chapters 3-4, 3-5 and 3-8 as the statement of what remains.

The thread begins with order. Theorem~50.1 reads the order coordinate as a count of one-steps, so that what the literature calls a coordinate along a worldline appears here as a count, and Theorem~62.1, at the end of the Part, closes the circle by showing that the group moves preserve the order parameter while the transfers raise it. The count is a count and never a flow: it is raised by one at every arrow step of the roster of the construction and at every depth step a read takes at its own star, and by no census move, so it is the translation of the four-coordinate lattice of counts along its order axis, with the dynamics the arrow operator of the construction, acyclic at all sixty-four networks, Theorem~50.1. Between these two results stand the chapters on the groups. Theorem~53.1 reads the rotation group as the automorphisms of the derived directions, Theorem~54.1 reads the Lorentz group as the isometries of the Hermitian split, and Theorem~56.1, an open result, places the boost at the one unit.

Chapter 3-2 is where a reader can see the whole method of the Part in one place. The rotation group of the theory is the stabilizer of the Hermitian unit inside the isometries of the Hermitian split of the trace, a Lie algebra of dimension three spanned by the ad matrices of the three derived directions with the brackets of the alternating symbol, Theorem~53.1. Its connected group is the set of exponentials at exact rational points, lifted to the qubit end as conjugation by its unitaries, and its finite lift is the dihedral image of order eight inside the twenty-four lattice rotations, Theorem~53.1. Nothing continuous is assumed: the directions are finite in number, the brackets are computed, and the group is what those brackets generate. Chapter 3-3 adjoins the boosts. The isometry algebra of the Hermitian split at the qutrit-qubit link, the plain-trace form of signature one and three at the Hermitian unit and the three derived directions, is six-dimensional, spanned by the ad matrices of the directions, which fix the unit, and by three boosts generated from the hyperbolic generator by the rotation brackets, and at every structure constant it equals the Lorentz algebra of the literature by name, with the Killing form of signature three and three, Theorem~54.1. The signature is therefore an output of the split and not a choice, which is the one fact of this Part a reader of relativity will want to check first.

Chapter 3-4 is an open result, and the introduction states its three faces plainly so that the reader does not mistake the boost for a theorem. No established rule mixes the order register with the three directions, the nine rules the theory carries on the five coordinates carrying a zero block between them, so the relative sign of the two coefficients of a boost is fixed by no rule, Theorem~56.1. The two ticks are exact and their squared ratio is two at the standing unit and one at two admitted alternates, so the integer maps preserving the Lorentzian form are the Pell family at the standing unit and the sign maps alone at the alternates, and the boost read from the tick ratio is unit-dependent, the one unit being the missing object, Theorem~56.1. What the construction has is the algebra of chapter 3-3 with its boosts as generators; what it does not have is the rule that would fix which integer map the dynamics selects, and the box names that absence as the parting.

Chapter 3-5 reads spin from the open-end label in Theorem~57.1, an open result whose parting the box names. The twice-spin read from the open qubit ends alone reaches the name at six of the nine bodies and one of the five carriers, the up quark first, Theorem~57.1. The composition of the doublet objects by the symmetric power per passage reads one doublet at every body and the one at no carrier, and the standing spin reading at the eight networks of the up-type quarks and the carriers equals the census doublet content at seven of the eight, the one parting re-set at the attach site, Theorem~57.1. The reader should notice the pattern that this open result shares with the boost open result: the construction reaches the name at most of the objects and parts at a nameable few, and the box is written so that the few are listed rather than averaged away.

Chapter 3-6 reads parity as the improper element in Theorem~59.1, with the discrete laws as the Weyl reflections at the cap in Open result~60.1 beside it as an open result that reaches forward to the gate set of Part XV. The lattice group of the three directions has forty-eight signed permutation matrices, half of them improper, while the identification group at the qubit end acts through eight proper ones and no identification unitary induces an improper element, so parity is not a symmetry of the records, Theorem~59.1. The open result beside it is the finite version of the discrete laws: at every one of the seventeen tabulated capped sinks the cap algebra of every link is the Weyl family of the act on the pair, its sixteen reflections typed at exact matrices over the ring of the twenty-fourth root, all sixteen linear at the (2,2) type, eight linear and eight antilinear at the (2,3) and (2,4) types, and four linear, four antilinear and eight absent at the three remaining types, with every reversal of the theory an arrow reversal at its support, Open result~60.1. Chapter 3-7 then proves the theorem the Part opened with. Over every edge of the standing move tables, the order parameter, the product over the links of a web of the dimension of the common support, is preserved by the group moves, the step and the holonomy, a gate changing no dimension, and raised by every transfer, the attach, the cut, the hub-forming move, the hub steps and the joins, at every read record edge with none non-rising, Theorem~62.1. Chapter 3-8 is counted empty, which the table titled The chapter counted empty (in chapter~\ref{mod:II-10-decoherence-and-the-open-system}, among the instruments) records: the chapter names the causal structure and general covariance, and the nearest established statements are the bound of one link per instance of chapter 3-1 and the lattice Poincare group of chapter 3-3.

In the language of the literature, this Part names the rotation group, the Lorentz group, the signature, the boost, spin, helicity, parity, charge conjugation, the reversal of order and their product, and general covariance. The correspondence is by name at each of them. What the boxes establish is the finite structure the names are attached to, the automorphisms of a finite set of directions and the isometries of a finite split. A physicist who reads the Lorentz algebra of chapter 3-3 as the Lorentz algebra is reading a correspondence that the box types as exact at the structure constants and as a name at the group, and the difference is the subject of the boost open result.

The dependency lines of the Part point to chapter 1-7 and to one another. The Lorentz group rests on the rotation group, the boost rests on the order coordinate and the Lorentz group, and parity rests on the rotation group and on spin. The order theorem of chapter 3-7 rests on the move tables of Part I and on the order coordinate of chapter 3-1, so that the first and the last chapter of the Part are one argument read at two ends. A reader following the derivation graph will find that the Parts on gauge structure and electrodynamics cite the rotation group and spin by number, that the discrete laws return at the gate set of Part XV, and that the gravity Part returns to the order coordinate as the count it reads the geometry against.
\section{The order coordinate and its theorem}\label{mod:III-1-order-coordinate-and-its-theorem}
\chdesig{4-1}{3-1}{III-1-order-coordinate-and-its-theorem}{III-1-D543820-w-tkud}
\label{narr:open:3-1}The order of events comes first. Theorem~50.1 reads the order coordinate as a count of one-steps, resting on the table of nouns of chapter 1-7.
\begin{chconditional}[The order coordinate as a count of one-steps]\label{box:III-1-order-coordinate-and-its-theorem:1}
This result is conditional on the axiom text and on the results of the chapters named as its premises in the derivation below, each at its own standing.
The order coordinate the compendium borrowed its word for is the count of one-steps of the order register, a count and never a flow: it is raised by one at every arrow step of the roster of the construction and at every depth step taken by a read at its own star, and by no census move, so it is the translation of the four-coordinate lattice of counts along its order axis; the dynamics is the arrow operator of the construction, acyclic at all sixty-four networks, and the depth step is posed and derived by no result.
\end{chconditional}
\chstanding{a conditional theorem}
\chdeps{Definitions 1.1 to 1.6}{Theorem 51.1, Open result 52.1, Theorem 56.1 and Open result 56.2, Theorem 62.1, Theorem 124.1 and Theorem 147.1}
The order step is an element of the one-dimensional centre of the frame algebra of the qutrit-qubit link; under the Killing form it is null, and under every trace form the ratio of the squared order step to the squared direction step is the exact rational four, so the standing tick ratio is one quarter, at which the integer maps preserving the Lorentzian pairing of the two ticks are the four sign maps alone.
Over every instance of every move of the fixed-row set the increment of the sign-ordered depth is zero, so the depth ticks under no move and is stepped by a read at its own star alone.
Over the full move census no instance of any move changes the washed register or any sign record at a read more than one link from the site of the move, so the influence of an instance reaches at most one link, a statement on the cells of the construction and no rate.
The register tower (a tower of Hilbert spaces of growing dimension), the ladder of registers of one pair type, gives the continuum as its limit: the readout differences along the tower shrink at every carrier of the construction and the limit read is fixed with no freedom by the tower, censused at the five carriers at rungs zero to twelve at two engines.
\begin{chbridge}
In the standard account the coordinate $t$ of Minkowski space is a continuous real parameter with the invariant interval $t^2 - x_1^2 - x_2^2 - x_3^2$ \cite{Minkowski1908}, the dynamics of a closed quantum system a one-parameter unitary group generated by a self-adjoint Hamiltonian \cite{Stone1932}, and a lattice theory reaches the continuum as the lattice spacing goes to zero \cite{Wilson1974}. The correspondence is by name: the compendium supplies an isometry algebra of the Hermitian split at the Lorentz brackets, a lattice point group, boosts from the order register and a transverse traceless part at the open read, every one a finite or algebraic object of this theory carrying the group name at chapter~\ref{mod:III-2-rotation-group}, chapter~\ref{mod:III-3-lorentz-group-and-the-signature}, chapter~\ref{mod:III-4-boosts-and-the-one-unit}, chapter~\ref{mod:III-5-spin-and-helicity}, the chapter titled Causal structure and general covariance (chapter~\ref{mod:III-8-causal-structure-and-general-covariance}, among the instruments) and chapter~\ref{mod:XIV-7-what-remains-the-ungraded-page}, and a boost as an element of the Lorentz group acting on a spacetime corresponds to these objects by name.
\end{chbridge}
\begin{chproof}[exact part and census part]
The premises of this derivation are the axiom and the six definitions of chapter~\ref{mod:FRAMEWORK} and the table of the chapter titled The Rosetta table in full (chapter~\ref{mod:I-7-rosetta-table-in-full}, among the instruments). The chain of premises runs back to the axiom through those chapters, and the longest such chain in this compendium is twelve chapters long.
The count is read off the arrow operator of the construction and its lift checked acyclic at every network; the centre of the frame algebra is shown one-dimensional and every trace form is shown to value the squared order step at four squared direction steps; the depth increment is computed at every instance from the commutation cells; the link distance of every changed register from the site of an instance is bounded over the census; and the tower is censused at the five carriers at two engines; each theorem is conditional on the results it names, among them the read-order and census-operator results, the ratio at the one unit and the depth monotonicity, and on the valence bound.
\textit{The exact part.}
From the arrow operator of the construction at every one of its sixty-four networks, by the lift along the order coordinate, we obtain an acyclic graph whose closing arcs are none, so the count of one-steps is a coordinate on every network and no cycle returns to its origin.
\begin{equation}
\text{closing arcs} = 0, \qquad \text{cyclic networks} = 0 \ \text{among } 64 .
\end{equation}
From the full move census, by the increment of the sign-ordered depth at every instance, we obtain the step census of the one-step: 2942 instances at one step, 609 at two and 148 at three among the 3699, the depth incremented by no census move and by a read at its own star alone.
\begin{equation}
\#\{\text{instances}\} = 2942 + 609 + 148 = 3699, \qquad \Delta_{\mathrm{depth}}(\text{move}) = 0 .
\end{equation}
From the one-dimensional center of the frame algebra of the qutrit-qubit link, by every trace form, we obtain the ratio of the squared order step to the squared direction step as the exact rational four, so the standing tick ratio is one quarter.
\begin{equation}
\frac{\langle s_{\mathrm{ord}}, s_{\mathrm{ord}} \rangle}{\langle s_{\mathrm{dir}}, s_{\mathrm{dir}} \rangle} = 4, \qquad \tau = \tfrac{1}{4} .
\end{equation}
\textit{The census part.}
By exhaustive enumeration of the 309 tabulated reads of the construction, the one-term action, the slot star alone, is the standing action at every read with no differing read, and under it every one of the six cluster cells is re-met at the own route with no differing cell. The arrow operator is re-met at all one hundred and ninety-two tables of its manifest and its lift checked acyclic at every network. The depth increment is computed at every instance of every move from the commutation cells, and the register tower is censused at the five carriers at rungs zero to twelve at two engines. The smallest case, the one-step census at the three step counts, is worked above.
The counter-tests alter the objects singly and read what moves. The two-term split moves four cells, the one-term action among them. The coordinate counter-test reads the same at 2942 instances and moves at 757, the step census re-met. The depth lift struck moves three cells. The planted text moves forty-seven. The two gate counter-tests are refused at the gate as pinned. The result was reproduced by a second engine built apart from the first, the two runs equal at every key, with the controls the same at both.
The machine-checked package named at the reproducibility appendix certifies the statements of this chapter under the definitions of this compendium and proves no theorem of the literature.
\chcert{3699}
\end{chproof}
\begin{figure}[h]\centering
\setlength{\unitlength}{1pt}
\begin{picture}(260,150)
\qbezier(86.0,75.0)(130.0,75.0)(174.0,75.0)
\put(130.0,83.0){\makebox(0,0){\scriptsize $(2,3)$}}
\put(77.0,75.0){\circle{14}}
\put(77.0,75.0){\makebox(0,0){\footnotesize $3$}}
\put(77.0,85.0){\makebox(0,0)[b]{\scriptsize $q$}}
\put(183.0,75.0){\circle{14}}
\put(183.0,75.0){\makebox(0,0){\footnotesize $2$}}
\put(183.0,85.0){\makebox(0,0)[b]{\scriptsize $b$}}
\end{picture}
\caption{The qutrit-qubit link whose frame algebra carries the order step at its one-dimensional center. The order coordinate is the count of one-steps of the order register along the link.}
\end{figure}
\section{One object at five readings}\label{mod:III-1-D543458-tkro-t1}
\begin{chconditional}[One object at five readings]\label{box:III-1-D543458-tkro-t1:1}
Conditional on the results this theorem rests on at their own standing, as its derivation names them.
The five open results of the order coordinate, the read-order open result, the reversal open result, the step-generator open result, the order-operator open result and the act open result, are five readings of one object, the serial (2,2) chain: the chain is built at nine of nine recorded webs, at the electron, the muon and the tau the sites a1, x1, a2, at the proton s1, x1, s4, at the neutron P.f1, P.x1, P.f2, at the photon word s3, x1, s2 and at the universe web B00.f1, B00.x1, B00.f2, the near end of dimension two at degree one or two, the inserted qubit and the far end of dimension two or three; its order is the three positions zero, one and two at one each, its reversal the far end, the inserted qubit and the near end, its step one site and one link added at the (2,2) genre of the first link, the site delta one and the link delta one at seven of seven, one link per tick as the one-link theorem bounds it, the two moves of the depth rule making the step at no chain and the step the first-parting kind of the depth-rule row at seven of seven; its office per tick is exact at seven of seven chains, $\tfrac{536569}{972} + \tfrac{25}{1458}\sqrt{2}$ at the electron, $\tfrac{847609}{12} + \tfrac{25}{18}\sqrt{2}$ at the muon, $\tfrac{3475947}{4} + \tfrac{25}{2}\sqrt{2}$ at the tau, zero at the proton and the neutron, $\tfrac{202201}{972} + \tfrac{25}{1458}\sqrt{2}$ at the photon word and $\tfrac{917593}{12} + \tfrac{25}{18}\sqrt{2}$ at the universe web. Against the five open results, each transcribed at its parting cell at the letter pinned, five of five: the passage of the read-order open result is defined at seven of seven and crosses the inserted qubit alone, the chain the first differing kind of that open result at seven of seven with the power register seven hundred twenty at four chains and three hundred eighty-four at the photon word; the reversal is defined and self-inverse at seven of seven with the reversed chain a recorded chain at none of seven, the electron the first; the one level of the energy open result equals the office per tick over the unit at the electron alone, the first differing the muon, equal to the first differing body of that open result; the positions equal the order spectrum at seven of seven while the transport generator at the chain link, of cycle type $1^{2}\,2^{1}$ with the exponents zero at three and twelve at one, parts from the three rungs by dimension at every chain; the transport at the chain link equals the two hundred fifty-five recorded (2,2) links of the act open result at seven of seven and parts from the exponential at one tick, a scalar, at every chain. By the rule that an open result is the chain when its object re-read at the chain is an exact cell at every recorded chain and its parting is re-met there, the readings that are the chain are five of five and the first that is not is none; the first face reads yes and the second face is unmet.
\end{chconditional}
\chstanding{conditional theorem}
\chdeps{Definitions 1.1 to 1.6, Theorem 50.1, Theorems 105.1 and 105.2, Open results 105.3 to 105.5 and Theorem 105.6 and Theorem 115.1}{no theorem beyond it at this release}
The five open results are transcribed at their parting cells and objects from the copies by hash; the serial (2,2) chain is built at every recorded web with its order, reversal and step; the office per tick is read at every chain; each open result's object is re-read at the chain by the substitution the item names and its parting re-met, and the readings that are the chain are counted.

\begin{chbridge}
A discrete dynamics is the repeated application of one step operator, the order coordinate of a chain is the index of its site, and the reversal, the generator and the single step are readings of the same translation \cite{Feynman1982,tHooft2016}. The derivation supplies the chain at every recorded web with the five open results re-read at it here, and the dynamics is named by correspondence.
\end{chbridge}
\begin{chproof}[census and exact]
The premises are the read-order open result, the reversal open result, the step-generator open result, the order-operator open result and the act open result (chapter 3-1, open results), the serial insertion row at the chains (chapter 10-6, an open result), the tick as the office of one serial link (chapter 10-6, an open result), the one-link theorem at the reach (chapter 10-6, a theorem) and the depth rule with its two moves (chapter 10-4, an open result). The chain back to the axiom runs through chapter 3-1 and the framework chapter. By exhaustive enumeration of the five open results at their parting cells, the nine recorded webs at their serial chains, the seven recorded chains at their offices per tick and the three thousand six hundred ninety-nine instances of the reversal open result re-met at the recorded instances, the property tested at each open result being whether its object re-read at the chain is an exact cell at every recorded chain and whether its parting is re-met there. From the near end, the inserted qubit and the far end at every web, by the positions, we obtain the chain and its three readings,
\begin{equation}
\text{order } (0, 1, 2) , \qquad \text{reversal } (2, 1, 0) , \qquad \text{step } \Delta_{\text{sites}} = 1 ,\ \Delta_{\text{links}} = 1 \ \text{at } (2,2) ,
\end{equation}
one link per tick at the reach one link per instance and zero per read, forced, and the recorded depth ticking under no move at the increment zero at eight of eight moves of the depth-step row; the hub-forming move adds two sites and the join keeps the site count, so the two moves of the rule make the step at no chain, the step the first-parting kind of the depth-rule row, an attach at the (2,2) genre, at seven of seven. From the office of one serial (2,2) link over the tick one, by the tick row, we obtain the office per tick at every chain, exact at seven of seven, the electron, the muon and the tau as the statement prints them with the proton and the neutron capped at zero. The five open results re-read at the chain, each with its object an exact cell at every chain and its parting re-met: the read-order open result, the passage defined at seven of seven crossing the inserted qubit alone so that the junction factor carries no part of the count, the count face zero and the power face two, the chain the first differing kind at seven of seven, the power register seven hundred twenty at the electron, the muon, the tau and the universe web equal to the first differing read of the open result and three hundred eighty-four at the photon word; the reversal open result, defined and self-inverse at seven of seven with the support two and the exponent image one, the reversed chain a recorded chain at none of seven, the first not a member the electron at the measure seven over seven, the open result's own parting re-met; the energy open result, the one level of the generator equal to the office per tick over the unit at one of seven, the electron alone, the first differing the muon equal to the first differing body of the open result and the bodies cell one of thirteen; the order-operator open result, the positions $0{:}1;1{:}1;2{:}1$ equal to the order spectrum at seven of seven, the transport generator at the chain link of cycle type $1^{2}\,2^{1}$ with the exponents zero at three and twelve at one over the ring of order twenty-four at the dimension four, against the three rungs equal at none of seven, the parting of dimension at every chain; the act open result, the transport at the chain link with two fixed points, two distinct eigenvalues and the characteristic polynomial $-1, 2, 0, -2, 1$ from the constant term equal to the two hundred fifty-five recorded (2,2) links at seven of seven, against the exponential at one tick, a scalar with one distinct eigenvalue, equal at none of seven, the kinds the phase alone at five and the identity at the two capped. So the readings that are the chain are five of five, which is the theorem. The smallest instance on the page is the electron chain a1, x1, a2 at the order zero, one, two, its reversal a2, x1, a1 and its step one site and one (2,2) link.

The counter-tests struck the inserted qubit from the chain, which left the step undefined at every web. A second put the hub-forming move in place of the attach as the step, which moved the site delta to two and the depth-rule reading off the chain. A third read the office per tick at the tick two, which halved every office with the energy open result reading unmoved at the electron alone, and a fourth struck the gate, which left every cell absent at both engines. Every control read the same at both engines, and the readings were reproduced at a second engine of a different construction, byte-identical at every run.
\chcert{3699}
\end{chproof}
\begin{figure}[h]\centering
\begin{adjustbox}{max width=\linewidth}\begin{tabular}{lll}\toprule
the open result & its object re-read at the chain & its parting re-met at the chain \\ \midrule
the read-order open result & the passage from the near end to the far end, defined at seven of seven & the chain the first differing kind at seven of seven \\
the reversal open result & the far end, the inserted qubit, the near end, self-inverse at seven of seven & a recorded chain at none of seven, the electron first \\
the step-generator open result & the one level against the office per tick, equal at the electron alone & the first differing the muon, equal to the open result \\
the order-operator open result & the positions zero, one, two, equal to the order spectrum at seven of seven & the transport generator parting by dimension at seven of seven \\
the act open result & the transport at the chain link, equal to the recorded links at seven of seven & the exponential at one tick a scalar, parting at seven of seven \\ \bottomrule
\end{tabular}\end{adjustbox}
\caption{The five open results of the order coordinate re-read at the serial chain, each object an exact cell at every recorded chain and each parting re-met there.}
\end{figure}
\section{The office per tick over the dressing squared}\label{mod:III-1-D543820-w-tkud}
\begin{chwall}[The office per tick over the dressing squared]\label{box:III-1-D543820-w-tkud:1}
The quotient of the office of the inserted serial (2,2) link over the square of the chain dressing is not one element at every recorded chain, the first parting the muon against the electron: the office of the inserted link at the adopted reading of the tick open result, re-met as the office at the first step less the office at the zeroth at nine of nine chains, is $\tfrac{536569}{972} + \tfrac{25}{1458}\sqrt{2}$ at the electron, $\tfrac{847609}{12} + \tfrac{25}{18}\sqrt{2}$ at the muon, $\tfrac{3475947}{4} + \tfrac{25}{2}\sqrt{2}$ at the tau, zero at the proton and the neutron, $\tfrac{202201}{972} + \tfrac{25}{1458}\sqrt{2}$ at the photon word and $\tfrac{917593}{12} + \tfrac{25}{18}\sqrt{2}$ at the universe web; the dressing of every chain is the two-point function of the six-dressings row at the genre of its recorded link, $\tfrac{5}{384} + \tfrac{5}{384}\sqrt{2}$ at the (2,2) link of the electron, the muon, the tau, the neutron and the universe web and $\tfrac{79}{31104} + \tfrac{1}{648}\sqrt{2} + \tfrac{7}{5184}\sqrt{3} + \tfrac{1}{10368}\sqrt{6}$ at the (2,3) link of the proton and the photon word, re-met at the Green-function cell at separation one at nine of nine; the quotient is computed at forty-five of forty-five cells, nine chains at five readings, and at the adopted reading it is $\tfrac{19779260416}{2025} - \tfrac{162791424}{25}\sqrt{2}$ at the electron, $\tfrac{31245438976}{25} - \tfrac{20830224384}{25}\sqrt{2}$ at the muon, $\tfrac{384404557824}{25} - \tfrac{256269090816}{25}\sqrt{2}$ at the tau, zero at the proton and the neutron and $\tfrac{33825329152}{25} - \tfrac{22550151168}{25}\sqrt{2}$ at the universe web, six distinct cells at the seven recorded chains. The first face, one element at every chain, reads no, and the second face reads yes with the first parting the muon against the electron; at the reading where the link-alone cell is taken over the chain dressing of the tick open result squared, the five uncapped chains stand at one element $\tfrac{15577}{972} + \tfrac{25}{1458}\sqrt{2}$ and the two capped at zero, a reading printed beside.
\end{chwall}
\chstanding{open result}
\chdeps{Definitions 1.1 to 1.6, Theorem 50.1, Open result 92.1 and Theorem 92.2, Theorem 115.1 and Theorems 185.1 to 185.12}{no theorem beyond it at this release}
The office of the inserted serial link is located at the nine chains by locator and re-met as a difference of offices; the recorded link of every chain is located with its genre and its dressing is read at the six-dressings row and re-met at the Green-function cell; the quotient of the office over the dressing squared is computed at five readings by the conjugate-norm inverse of the field and compared across the chains.

\begin{chbridge}
The renormalised propagator of a field is the bare propagator divided by the square of the wave-function renormalisation, one factor of the field strength at each end of the line \cite{Dyson1949,LSZ1955}. The derivation supplies the quotient at every recorded chain here, and the renormalisation is named by correspondence.
\end{chbridge}
\begin{chproof}[finite and exact]
The premises are the tick as the office of one serial link (chapter 10-6, an open result), the Green-function row at the dressing (chapter 10-2, an open result), the six dressings at the six genres (chapter 8-1, an open result), the serial insertion row at the chains (chapter 10-6, an open result) and the Ward-functional row at the (2,3) dressing (chapter 4-8, an open result). The chain back to the axiom runs through chapter 10-0 and the framework chapter. The procedure. The input is the nine chains, the seven recorded and the two beside, each with its office of the inserted serial (2,2) link at the adopted reading of the tick open result, located at nine of nine, and its recorded link with its genre, a (2,2) link at the electron, the muon, the tau, the neutron and the universe web and a (2,3) link at the proton and the photon word, the capped chains the item names. The enumeration rule reads the dressing of every chain as the two-point function of the six-dressings row at that genre, computes the quotient at five readings, the adopted reading the office over the dressing squared, the link-alone cell over the dressing squared, the office over the chain dressing of the tick open result squared, the link-alone cell over the chain dressing squared and the link-alone cell over the two dressings multiplied and squared, and compares the cells across the chains; the test at the adopted reading is whether the quotient is one element at every recorded chain. From the office of the electron chain and the (2,2) dressing, by the conjugate-norm inverse of the field, we obtain
\begin{equation}
\frac{\tfrac{536569}{972} + \tfrac{25}{1458}\sqrt{2}}{\left(\tfrac{5}{384} + \tfrac{5}{384}\sqrt{2}\right)^{2}} = \frac{19779260416}{2025} - \frac{162791424}{25}\sqrt{2} ,
\end{equation}
and alike at the muon $\tfrac{31245438976}{25} - \tfrac{20830224384}{25}\sqrt{2}$, the identity quotient multiplied by the dressing squared equal to the office at forty-five of forty-five cells. The counts by kind: the office re-met as the office at the first step less the office at the zeroth at nine of nine chains and the link-alone cell, the act of the inserted link squared plus the closure increment, re-met at nine of nine with the electron and the photon word at one cell; the (2,2) dressing equal at the Lean spelling, the zero-displacement row, the one-link cell of the Green-function row and the unit cell, the (2,3) dressing equal at the Lean spelling, the Ward-functional row and the unit cell; the quotient at the adopted reading six distinct cells at seven recorded chains, one element at every chain no, the first parting the muon against the electron, $\tfrac{31245438976}{25} - \tfrac{20830224384}{25}\sqrt{2}$ against $\tfrac{19779260416}{2025} - \tfrac{162791424}{25}\sqrt{2}$; at the readings beside the one-element letter no at every one, with the five uncapped chains at one element and the two capped at zero where the link-alone cell is taken over the chain dressing squared, and the four (2,2)-link uncapped chains at one element where it is taken over the two dressings multiplied, the photon word parting at its (2,3) dressing. So the quotient parts at the muon first, which is the open result. The smallest instance on the page is the electron chain, the office over the dressing squared at the display above.

The counter-tests struck the dressing to one, which moved every quotient to the office itself with the partings standing. A second read the (2,3) dressing at every chain, which moved the five (2,2) chains and left the proton and the photon word. A third read the link-alone cell in place of the office, which moved the five uncapped chains to one element at the chain-dressing reading, and a fourth struck the gate, which left every cell absent at both engines. Every control read the same at both engines, and the forty-five cells were reproduced at a second engine, equal at every cell.
\chcert{45}
\end{chproof}
\begin{figure}[h]\centering
\begin{adjustbox}{max width=\linewidth}\begin{tabular}{llll}\toprule
chain & the office of the inserted link & the dressing & the quotient at the adopted reading \\ \midrule
the electron & $\tfrac{536569}{972} + \tfrac{25}{1458}\sqrt{2}$ & $\tfrac{5}{384} + \tfrac{5}{384}\sqrt{2}$ & $\tfrac{19779260416}{2025} - \tfrac{162791424}{25}\sqrt{2}$ \\
the muon & $\tfrac{847609}{12} + \tfrac{25}{18}\sqrt{2}$ & $\tfrac{5}{384} + \tfrac{5}{384}\sqrt{2}$ & $\tfrac{31245438976}{25} - \tfrac{20830224384}{25}\sqrt{2}$ \\
the proton & $0$ & the (2,3) dressing & $0$ \\
the universe web & $\tfrac{917593}{12} + \tfrac{25}{18}\sqrt{2}$ & $\tfrac{5}{384} + \tfrac{5}{384}\sqrt{2}$ & $\tfrac{33825329152}{25} - \tfrac{22550151168}{25}\sqrt{2}$ \\ \bottomrule
\end{tabular}\end{adjustbox}
\caption{The office of the inserted serial link, the dressing and their quotient at four of the nine chains, the first parting the muon against the electron.}
\end{figure}
\par\medskip\noindent{\itshape With order read as a count, the next chapter turns to the rotations.}\label{narr:close:3-1}\par
\section{The rotation group}\label{mod:III-2-rotation-group}
\chdesig{4-2}{3-2}{III-2-rotation-group}{III-2-rotation-group}
\label{narr:open:3-2}Theorem~53.1 reads the rotation group as the automorphisms of the derived directions, the group of the literature appearing as the symmetries of a finite set of directions the construction derives.
\begin{chconditional}[The rotation group as the automorphisms of the derived directions]\label{box:III-2-rotation-group:1}
This result is conditional on the axiom text and on the results of the chapters named as its premises in the derivation below, each at its own standing.
The rotation group of the theory is the stabiliser of the Hermitian unit inside the isometries of the Hermitian split of the trace, a Lie algebra of dimension three spanned by the ad matrices of the three derived directions with the brackets of the alternating symbol, its connected group the exponentials at exact rational points lifted to the qubit end as conjugation by its unitaries, the finite lift the dihedral image of order eight inside the twenty-four lattice rotations; every qubit end of the thirteen bodies of the construction is typed by it, properly at the near end of a dock and improperly at the far end of a qubit-qubit dock, so the dock qubit ends are spatial and the hub-leg copies internal.
\end{chconditional}
\chstanding{a conditional theorem}
\chdeps{Definitions 1.1 to 1.6}{Open result 17.1, Theorem 54.1, Theorem 57.1 and Open result 57.2, Theorem 59.1, Open result 123.1 and Theorem 124.1}
The identification group at the qubit end of a read link, thirty-two monomial matrices with fourth-root phases, acts on the three directions through its image, the dihedral group of order eight, a subgroup of index three in the twenty-four lattice rotations, with the kernel the four scalar phases; the directions carry $A_2 \oplus E$ and the pair space of the end $A_1 \oplus A_2 \oplus E$, so its commutant has dimension three; the readout of every body is unchanged under the group, an identity; the orbital label of a body is conserved at every interaction instance, conditional on the one-act clause; and a rotation by a generic angle lies outside the group, the continuum closed at the first result above.
Two open results stand here. One direction per pair from the shared support, the second clause of the read definition, is a uniqueness claim where the construction admits several readings, decided at one of four and superseded by the automorphism theorem. The full rotation is minus one at every qubit end and plus one on the directions, and the own-site sign of a body read as its phase at the spatial ends equals the standing swap sign at seventeen of the eighteen objects of the construction, parting at the tau, whose own site is a qubit-qubit port paired to a leaf the transport opens at neither orientation.
\begin{chbridge}
The rotation group of three-dimensional space is $SO(3)$, its universal cover $SU(2)$, with the Lie algebra spanned by three generators at the brackets $[T_a, T_b] = \epsilon_{abc} T_c$; a lattice admits a finite point group, the rotations of a square forming the dihedral group of order eight with the irreducible representations $A_1$, $A_2$, $B_1$, $B_2$ and $E$ \cite{Hamermesh1962}; and the representation of the rotation group on a quantum system is projective, lifted to $SU(2)$ \cite{Wigner1959}. The correspondence is by name: the compendium supplies a web of this theory read at its sinks, a doublet swap sign at the odd cycle, a Dirac form at an established generator and a finite record sum carrying the path-integral name at chapter~\ref{mod:II-1-state-space-from-the-act-outputs}, chapter~\ref{mod:II-5-indistinguishability-and-the-spin-statistics-sign}, chapter~\ref{mod:II-7-path-integral-as-the-finite-record-sum}, chapter~\ref{mod:III-5-spin-and-helicity}, the chapter titled The Dirac equation at the dock (chapter~\ref{mod:V-9-dirac-equation-at-the-dock}, among the instruments), chapter~\ref{mod:VI-2-w-and-the-z-and-their-masses}, chapter~\ref{mod:VI-2b-neutral-lepton-the-hub-unit-with-its-legs-open}, chapter~\ref{mod:VII-5-proton-to-electron-ratio-a-candidate}, chapter~\ref{mod:VIII-1-mass-unit-and-the-first-power}, chapter~\ref{mod:XIV-2-walls-census-in-physics-terms} and chapter~\ref{mod:XIV-7-what-remains-the-ungraded-page}, a finite re-matching group at the known clause, a representation census at the open ends and the sinks, an attach-site basis carrying the generation name and a sink class carrying the doublet name at chapter~\ref{mod:IV-1-gauge-invariance-as-the-known-clause}, chapter~\ref{mod:IV-3-representations-from-the-open-end-census}, chapter~\ref{mod:IV-4-cartan-weights-and-the-charge}, chapter~\ref{mod:IV-6-conservation-as-the-kernel-of-the-move-matrix}, chapter~\ref{mod:IV-7-noether-at-the-sinks-and-the-kernel-identification}, chapter~\ref{mod:IV-9-one-rule-of-a-read-and-the-word-table}, chapter~\ref{mod:VI-3-doublets-handedness-and-parity-violation}, chapter~\ref{mod:VII-7-quark-content-and-flavour}, chapter~\ref{mod:VIII-4-three-generations-as-attach-sites}, chapter~\ref{mod:XIV-7-what-remains-the-ungraded-page} and the Rosetta table of this theory nouns, a pair stabilizer, the cap records as a Weyl basis, the Weyl relation at the uncertainty bound, the Bell bound at the pair ket, the gate set at the Weyl criterion, the wash at a sub-web and the Born count as the dressed measure at chapter~\ref{mod:FRAMEWORK}, chapter~\ref{mod:I-1-axiom-and-the-six-definitions}, chapter~\ref{mod:I-3-records-reads-and-the-born-count-the-gain-lemma}, chapter~\ref{mod:I-4-moves-in-one-table}, chapter~\ref{mod:I-6-read-at-a-register-the-probe-and-the-measurement}, the chapter titled The Rosetta table in full (chapter~\ref{mod:I-7-rosetta-table-in-full}, among the instruments), chapter~\ref{mod:II-1-state-space-from-the-act-outputs}, chapter~\ref{mod:II-2-born-rule-and-unitarity}, chapter~\ref{mod:II-3-uncertainty-relation-from-the-weyl-relation}, chapter~\ref{mod:II-9-qm-corollaries-the-fifteen-entries}, the chapter titled Decoherence and the open system (chapter~\ref{mod:II-10-decoherence-and-the-open-system}, among the instruments), chapter~\ref{mod:III-6-p-c-t-and-cpt}, chapter~\ref{mod:III-7-order-theorem-group-moves-and-transfers}, chapter~\ref{mod:IV-2-gauge-group-as-a-corollary-of-one-construction}, chapter~\ref{mod:V-1-fine-structure-constant-at-the-one-link}, chapter~\ref{mod:VI-2-w-and-the-z-and-their-masses}, chapter~\ref{mod:VI-3-doublets-handedness-and-parity-violation}, chapter~\ref{mod:VIII-2-body-alone-office-rule}, chapter~\ref{mod:XI-4-order-parameter-at-the-stop-condition}, chapter~\ref{mod:XI-7-thermodynamic-remainder-and-the-arrow}, chapter~\ref{mod:XIII-1-closure-as-the-computation}, chapter~\ref{mod:XIII-2-moves-as-a-gate-set-and-the-weyl-criterion}, chapter~\ref{mod:XIII-4-distillation-the-read-as-the-maximal-distillation-and-the-distillation-theorem}, chapter~\ref{mod:XIII-5-cost-of-a-read-and-the-adiabatic-class}, chapter~\ref{mod:XIII-6-holonomy-and-the-record-transport}, chapter~\ref{mod:XIV-7-what-remains-the-ungraded-page}, and the tau as the spin-one-half quantum of a field with its mass and spin as Lorentz labels, the factor named as a Lie group acting on the fields of a Lagrangian with its representations and the object named as a theorem or a standard object of a Hilbert space with its tensor structure correspond to these objects by name.
\end{chbridge}
\begin{chproof}[census]
The premises of this derivation are the axiom and the six definitions of chapter~\ref{mod:FRAMEWORK} and the table of the chapter titled The Rosetta table in full (chapter~\ref{mod:I-7-rosetta-table-in-full}, among the instruments). The chain of premises runs back to the axiom through those chapters, and the longest such chain in this compendium is twelve chapters long.
Inside the isometries of the Hermitian split the stabiliser of the unit is shown the span of the ad matrices of the three derived directions, and the point group is generated by the swap and the diagonal phase, a quarter turn and a half turn, its characters decomposing the directions and the pair space. The automorphism theorem rests on the axiom text and on two named conditions, that the directions of a link are those of its shared support at the end read out and that a dock pairs at its standing orientation; the point-group theorem on the one-act clause and on the open results it rests on.
By exhaustive enumeration of the one hundred and two qubit ends of the bodies of the construction at their sites, links and pair types, the stabilizer of the Hermitian unit inside the isometries of the Hermitian split is the span of the ad matrices of the three derived directions, a Lie algebra of dimension three with the brackets of the alternating symbol, its connected group the exponentials at exact rational points lifted to the qubit end as conjugation by its unitaries. The identification group at the qubit end, thirty-two monomial matrices with fourth-root phases, acts on the three directions through its image, the dihedral group of order eight, of index three in the twenty-four lattice rotations with the kernel the four scalar phases, and every qubit end is typed by it, properly at the near end of a dock and improperly at the far end of a qubit-qubit dock. The smallest case is the point group at one qubit end, generated by the swap and the diagonal phase, whose eight elements decompose the three directions as the sign representation plus the two-dimensional one and the pair space as the trivial, the sign and the two-dimensional representations, so the commutant has dimension three and the index of the image is twenty-four over eight.
The counter-tests alter the objects singly and read what moves, the unit struck leaving no group and an isometry algebra of dimension three. The five counter-tests stand at their pinned letters, the thirteen controls read the same at both engines and a second engine built apart reproduced every cell.
The machine-checked package named at the reproducibility appendix certifies the statements of this chapter under the definitions of this compendium and proves no theorem of the literature.
\chcert{102}
\end{chproof}
\par\medskip\noindent{\itshape The rotations fixed, the next chapter joins them to the order coordinate in the Lorentz group.}\label{narr:close:3-2}\par
\section{The Lorentz group and the signature}\label{mod:III-3-lorentz-group-and-the-signature}
\chdesig{4-3}{3-3}{III-3-lorentz-group-and-the-signature}{III-3-D543402-emis-t1}
\label{narr:open:3-3}Theorem~54.1 reads the Lorentz group as the isometries of the Hermitian split, and the signature with it, resting on the rotation group of the preceding chapter.
\begin{chconditional}[The Lorentz group as the isometries of the Hermitian split]\label{box:III-3-lorentz-group-and-the-signature:1}
This result is conditional on the axiom text and on the results of the chapters named as its premises in the derivation below, each at its own standing.
The isometry algebra of the Hermitian split at the qutrit-qubit link, the plain-trace form of signature one and three at the Hermitian unit and the three derived directions, is six-dimensional, spanned by the ad matrices of the directions, which fix the unit, and by three boosts generated from the hyperbolic generator by the rotation brackets; at every structure constant it equals the Lorentz algebra of the literature by name, with the Killing form of signature three and three, and the six lift to the qubit end as three anti-Hermitian and three Hermitian generators.
\end{chconditional}
\chstanding{a conditional theorem}
\chdeps{Definitions 1.1 to 1.6 and Theorem 53.1}{Theorem 55.1, Theorem 56.1 and Open result 56.2, Theorem 58.1, Open result 88.1, Open result 123.1 and Open result 236.1}
The Hermitian order unit, minus $i$ multiplied by the centre of the frame algebra, lies outside the anti-Hermitian algebra, and under the plain trace the real span of the unit and the three direction steps carries an exact Gram matrix, the unit squaring to the link dimension six, every direction step to minus three halves and every interference term (the cross term) zero, so the Gram matrix is three halves multiplied by the Lorentzian pairing of the lattice of ticks, the signature one plus three, the same split holding at every one of the six pair types.
In the four coordinates of the lattice of counts the integer maps preserving the Lorentzian form contain an element of infinite order at every one of the eight squared tick ratios the tick results admit, while the Euclidean pairing admits a finite group at every ratio, so the signature by existence reads Lorentzian at all eight ratios, the one unit open at chapter~\ref{mod:III-4-boosts-and-the-one-unit}.
With the four lattice translations, the point group of chapter~\ref{mod:III-2-rotation-group} and the first boost of the Pell family, the lattice Poincar\'e group of the theory exists and closes as a semidirect product at every placement of the boost and at all four admitted units, every element an integer affine map of the lattice and the tabulated readouts invariant under it.
\begin{chbridge}
Minkowski space carries the quadratic form $t^2 - x_1^2 - x_2^2 - x_3^2$ of signature one plus three, the Lorentz group $O(3,1)$ its group of linear isometries with a Lie algebra of dimension six, three rotations and three boosts, and the Poincar\'e group the semidirect product of the translations with it \cite{Minkowski1908,Weinberg1995}; on the integer lattice the maps preserving an indefinite form such as $t^2 - 2x^2$ are the infinitely many solutions of a Pell equation, while the maps preserving a positive definite form are finitely many \cite{HardyWright1979}. The correspondence is by name: the compendium supplies the Rosetta table of this theory nouns, a pair stabilizer, the cap records as a Weyl basis, the Weyl relation at the uncertainty bound, the Bell bound at the pair ket, the gate set at the Weyl criterion, the wash at a sub-web and the Born count as the dressed measure at chapter~\ref{mod:FRAMEWORK}, chapter~\ref{mod:I-1-axiom-and-the-six-definitions}, chapter~\ref{mod:I-3-records-reads-and-the-born-count-the-gain-lemma}, chapter~\ref{mod:I-4-moves-in-one-table}, chapter~\ref{mod:I-6-read-at-a-register-the-probe-and-the-measurement}, the chapter titled The Rosetta table in full (chapter~\ref{mod:I-7-rosetta-table-in-full}, among the instruments), chapter~\ref{mod:II-1-state-space-from-the-act-outputs}, chapter~\ref{mod:II-2-born-rule-and-unitarity}, chapter~\ref{mod:II-3-uncertainty-relation-from-the-weyl-relation}, chapter~\ref{mod:II-9-qm-corollaries-the-fifteen-entries}, the chapter titled Decoherence and the open system (chapter~\ref{mod:II-10-decoherence-and-the-open-system}, among the instruments), chapter~\ref{mod:III-6-p-c-t-and-cpt}, chapter~\ref{mod:III-7-order-theorem-group-moves-and-transfers}, chapter~\ref{mod:IV-2-gauge-group-as-a-corollary-of-one-construction}, chapter~\ref{mod:V-1-fine-structure-constant-at-the-one-link}, chapter~\ref{mod:VI-2-w-and-the-z-and-their-masses}, chapter~\ref{mod:VI-3-doublets-handedness-and-parity-violation}, chapter~\ref{mod:VIII-2-body-alone-office-rule}, chapter~\ref{mod:XI-4-order-parameter-at-the-stop-condition}, chapter~\ref{mod:XI-7-thermodynamic-remainder-and-the-arrow}, chapter~\ref{mod:XIII-1-closure-as-the-computation}, chapter~\ref{mod:XIII-2-moves-as-a-gate-set-and-the-weyl-criterion}, chapter~\ref{mod:XIII-4-distillation-the-read-as-the-maximal-distillation-and-the-distillation-theorem}, chapter~\ref{mod:XIII-5-cost-of-a-read-and-the-adiabatic-class}, chapter~\ref{mod:XIII-6-holonomy-and-the-record-transport}, chapter~\ref{mod:XIV-7-what-remains-the-ungraded-page}, and the object named as a theorem or a standard object of a Hilbert space with its tensor structure corresponds to these objects by name.
\end{chbridge}
\begin{chproof}[exact part and census part]
The premises of this derivation are the axiom and the six definitions of chapter~\ref{mod:FRAMEWORK}, the table of the chapter titled The Rosetta table in full (chapter~\ref{mod:I-7-rosetta-table-in-full}, among the instruments) and the theorem of chapter~\ref{mod:III-2-rotation-group}. The chain of premises runs back to the axiom through those chapters, and the longest such chain in this compendium is twelve chapters long.
From the split alone the stabiliser of the unit is the rotation algebra of chapter~\ref{mod:III-2-rotation-group} and the hyperbolic generator solved from the Gram matrix generates three boosts under the rotation brackets; the Gram matrix is computed cell for cell at every pair type; a solution of the Pell equation gives a map of infinite order at every admitted ratio, the Euclidean finiteness following from the finitely many lattice points at a given value of a positive definite form; and the translations, the point group and the first boost are written as integer affine maps with their closure checked. The algebra theorem rests on the axiom text and on the rotation-group result at its two conditions; the others on the readings of the construction and on the open results they rest on.
\textit{The exact part.}
From the plain-trace form at the qutrit-qubit link restricted to the real span of the Hermitian unit and the three direction steps, by the Gram matrix computed cell for cell, we obtain the unit squaring to the link dimension six, every direction step to minus three halves and every cross term zero, so the Gram matrix is three halves multiplied by the Lorentzian pairing of the lattice of ticks, the signature one plus three.
\begin{equation}
G = \begin{pmatrix} 6 & 0 & 0 & 0 \\ 0 & -\tfrac{3}{2} & 0 & 0 \\ 0 & 0 & -\tfrac{3}{2} & 0 \\ 0 & 0 & 0 & -\tfrac{3}{2} \end{pmatrix} = \tfrac{3}{2}\,\operatorname{diag}(4, -1, -1, -1) .
\end{equation}
From the three ad matrices of the directions, which fix the unit, and the hyperbolic generator solved from the Gram matrix, by the rotation brackets, we obtain three further generators, every one a boost by its action on the unit, and the six close with exact structure constants, the first two rotations bracket to the third, the brackets of a rotation with a boost rotate the boost, and two boosts bracket to minus a rotation, the Jacobi identity holding and the Killing form of signature three and three.
\begin{equation}
[J_{1}, J_{2}] = J_{3}, \qquad [J_{i}, K_{j}] = \epsilon_{ijk} K_{k}, \qquad [K_{1}, K_{2}] = -J_{3}, \qquad \operatorname{sig}\,\kappa = (3, 3) .
\end{equation}
\textit{The census part.}
By exhaustive enumeration of the six pair types of the construction, the Gram matrix of the Hermitian split is computed cell for cell at every type and the same split of signature one plus three holds at every one. The isometry algebra of the split at the qutrit-qubit link is six-dimensional with every structure constant exact and equal by name to the Lorentz algebra of the literature, the six generators lifting to the qubit end as three anti-Hermitian and three Hermitian generators, and the gate of the derivation finds all thirty-four of its clauses. The smallest case, the Gram matrix and the brackets at the qutrit-qubit link, is worked above.
The counter-tests alter the objects singly and read what moves. With the unit struck the six-dimensional algebra moves to the three-dimensional isometry algebra of the directions alone at the signature zero and three, with no plane of the unit. The five counter-tests stand at their pinned letters and the thirteen controls read the same at both engines. The result was reproduced by a second engine built apart from the first, the two runs equal at every one of the cells of both.
The machine-checked package named at the reproducibility appendix certifies the statements of this chapter under the definitions of this compendium and proves no theorem of the literature.
\chcert{6}
\end{chproof}
\begin{figure}[h]\centering
\setlength{\unitlength}{1pt}
\begin{picture}(260,150)
\qbezier(86.0,75.0)(130.0,75.0)(174.0,75.0)
\put(130.0,83.0){\makebox(0,0){\scriptsize $(2,3)$}}
\put(77.0,75.0){\circle{14}}
\put(77.0,75.0){\makebox(0,0){\footnotesize $3$}}
\put(77.0,85.0){\makebox(0,0)[b]{\scriptsize $q$}}
\put(183.0,75.0){\circle{14}}
\put(183.0,75.0){\makebox(0,0){\footnotesize $2$}}
\put(183.0,85.0){\makebox(0,0)[b]{\scriptsize $b$}}
\end{picture}
\caption{The qutrit-qubit link at which the Hermitian split of the plain trace is taken, the Hermitian unit and the three derived directions span the real four-dimensional space whose Gram matrix carries the signature one plus three.}
\end{figure}
\section{One against three at every step of the chain}\label{mod:III-3-D543402-emis-t1}
\begin{chtheorem}[One against three at every step of the chain]\label{box:III-3-D543402-emis-t1:1}
At every chain and step carrying a (2,2) link, the Gram of the propagation reads one against three, the boost of the chain is its isometry with the mixing block nonzero, and the one-unit rotation of the boost open result is the isometry of a definite plane and of no restriction of that Gram: along the serial (2,2) chain of every recorded web of the serial insertion row at the steps zero to three, thirty-four of thirty-six chain states walked at one hundred sixteen steps over eighty-two serial (2,2) links, the carried unit and the three carried directions give the Gram $4,0,0,0;0,-1,0,0;0,0,-1,0;0,0,0,-1$ on the link space at both ends, equal to the (2,2)-pair row, its signature one against three at one hundred sixteen of one hundred sixteen steps equal to the pair-end signature, definite at none and other at none; the boost of the chain at every step is the plane $4,0;0,-1$ with the generator $0,\tfrac{1}{2};2,0$ at the residual zero, an isometry at one hundred sixteen of one hundred sixteen; the one-unit rotation of the open result at every step is the order-step plane $1,0;0,\tfrac{1}{4}$, definite, with the generator $0,-\tfrac{1}{2};2,0$ at the residual zero on the Euclidean pairing and $0,-4;-4,0$ on the Gram of the propagation restricted, Euclidean at the unit and Lorentzian at the chain; the orientation of the carried triple is preserved at the even steps and reversed at the odd, fifty of one hundred sixteen, with the structure constants at the re-oriented basis equal to the brackets table at ninety of ninety at one hundred sixteen of one hundred sixteen steps; the boost open result is transcribed at fifteen of fifteen cells and re-derived at the (2,3) link frame at ten of ten. The first face reads yes and the second face, the first chain and step parting named, reads no, the first parting none.
\end{chtheorem}
\chstanding{theorem}
\chdeps{Definitions 1.1 to 1.6, Theorem 54.1, Open result 92.1 and Theorem 92.2 and Theorem 115.1}{no theorem beyond it at this release}
The boost open result is transcribed from its copies by hash and re-derived at the (2,3) link frame; the serial (2,2) chain of every recorded web is propagated step by step with the carried unit and directions, the Gram of the propagation is read with its signature by the pivot signs and by the leading principal minors, the boost of the chain and the one-unit rotation of the open result are tested as isometries at the two pairings at every step, and the orientation of the carried triple is read against the brackets table.
\begin{figure}[htbp]\centering
\setlength{\unitlength}{1mm}
\begin{picture}(128,74)(0,0)
\put(14,62){\line(1,0){24}}\put(26,63){\makebox(0,0)[b]{\tiny (2,2)}}
\put(12,62){\circle{4}}\put(12,62){\makebox(0,0){\tiny p}}\put(12,59){\makebox(0,0)[t]{\tiny a1}}
\put(40,62){\circle{4}}\put(40,62){\makebox(0,0){\tiny r}}\put(40,59){\makebox(0,0)[t]{\tiny a2}}
\put(10,56){\makebox(0,0)[l]{\tiny the step k0: the Gram diagonal 4, -1, -1, -1, the signature 1 against 3}}
\put(14,46){\line(1,0){24}}\put(26,47){\makebox(0,0)[b]{\tiny (2,2)}}
\put(42,46){\line(1,0){24}}\put(54,47){\makebox(0,0)[b]{\tiny (2,2)}}
\put(12,46){\circle{4}}\put(12,46){\makebox(0,0){\tiny p}}\put(12,43){\makebox(0,0)[t]{\tiny a1}}
\put(40,46){\circle{4}}\put(40,46){\makebox(0,0){\tiny r}}\put(40,43){\makebox(0,0)[t]{\tiny x1}}
\put(68,46){\circle{4}}\put(68,46){\makebox(0,0){\tiny p}}\put(68,43){\makebox(0,0)[t]{\tiny a2}}
\put(10,40){\makebox(0,0)[l]{\tiny the step k1: the Gram diagonal 4, -1, -1, -1, the signature 1 against 3}}
\put(14,30){\line(1,0){24}}\put(26,31){\makebox(0,0)[b]{\tiny (2,2)}}
\put(42,30){\line(1,0){24}}\put(54,31){\makebox(0,0)[b]{\tiny (2,2)}}
\put(70,30){\line(1,0){24}}\put(82,31){\makebox(0,0)[b]{\tiny (2,2)}}
\put(12,30){\circle{4}}\put(12,30){\makebox(0,0){\tiny p}}\put(12,27){\makebox(0,0)[t]{\tiny a1}}
\put(40,30){\circle{4}}\put(40,30){\makebox(0,0){\tiny r}}\put(40,27){\makebox(0,0)[t]{\tiny x1}}
\put(68,30){\circle{4}}\put(68,30){\makebox(0,0){\tiny p}}\put(68,27){\makebox(0,0)[t]{\tiny x2}}
\put(96,30){\circle{4}}\put(96,30){\makebox(0,0){\tiny r}}\put(96,27){\makebox(0,0)[t]{\tiny a2}}
\put(10,24){\makebox(0,0)[l]{\tiny the step k2: the Gram diagonal 4, -1, -1, -1, the signature 1 against 3}}
\put(14,14){\line(1,0){24}}\put(26,15){\makebox(0,0)[b]{\tiny (2,2)}}
\put(42,14){\line(1,0){24}}\put(54,15){\makebox(0,0)[b]{\tiny (2,2)}}
\put(70,14){\line(1,0){24}}\put(82,15){\makebox(0,0)[b]{\tiny (2,2)}}
\put(98,14){\line(1,0){24}}\put(110,15){\makebox(0,0)[b]{\tiny (2,2)}}
\put(12,14){\circle{4}}\put(12,14){\makebox(0,0){\tiny p}}\put(12,11){\makebox(0,0)[t]{\tiny a1}}
\put(40,14){\circle{4}}\put(40,14){\makebox(0,0){\tiny r}}\put(40,11){\makebox(0,0)[t]{\tiny x1}}
\put(68,14){\circle{4}}\put(68,14){\makebox(0,0){\tiny p}}\put(68,11){\makebox(0,0)[t]{\tiny x2}}
\put(96,14){\circle{4}}\put(96,14){\makebox(0,0){\tiny r}}\put(96,11){\makebox(0,0)[t]{\tiny x3}}
\put(124,14){\circle{4}}\put(124,14){\makebox(0,0){\tiny p}}\put(124,11){\makebox(0,0)[t]{\tiny a2}}
\put(10,8){\makebox(0,0)[l]{\tiny the step k3: the Gram diagonal 4, -1, -1, -1, the signature 1 against 3}}
\put(4,70){\footnotesize the serial (2,2) chain of the electron web at four states; p the orientation preserved, r reversed}
\end{picture}
\caption{The propagation along the serial (2,2) chain of the electron web at the steps zero to three: the Gram of the propagation reads 4, -1, -1, -1 on its diagonal at every step, the signature one against three, and the orientation of the carried triple is preserved at the even steps and reversed at the odd steps. The sites carry the orientation letter, p preserved and r reversed, and every link the cell (2,2).}\label{fig:III-3-D543402-serial-chain}
\end{figure}
\begin{figure}[htbp]\centering
\setlength{\unitlength}{1mm}
\begin{picture}(160,64)(0,0)
\qbezier(80,24)(80,42)(80,60)
\qbezier(44,24)(80,24)(116,24)
\qbezier(80,24)(71,42)(62,60)
\qbezier(80,24)(89,42)(98,60)
\put(82,56){\makebox(36,5){\tiny the propagated unit, the entry 4}}
\put(116,20){\makebox(44,5){\tiny a carried direction, the entry -1}}
\put(98,60){\makebox(40,5){\tiny the null directions of the plane 4,0;0,-1}}
\put(22,60){\makebox(40,5){\tiny the boost generator 0,1/2;2,0}}
\put(0,50){\makebox(44,5){\tiny the signature 1,3, the ratio -4}}
\put(12,6){\framebox(12,7){\scriptsize a1}}
\put(8,1){\makebox(20,4){\tiny preserved}}
\qbezier(24,9)(34,9)(44,9)
\put(24,10){\makebox(20,4){\tiny (2,2)}}
\put(44,6){\framebox(12,7){\scriptsize x1}}
\put(40,1){\makebox(20,4){\tiny reversed}}
\qbezier(56,9)(66,9)(76,9)
\put(56,10){\makebox(20,4){\tiny (2,2)}}
\put(76,6){\framebox(12,7){\scriptsize x2}}
\put(72,1){\makebox(20,4){\tiny preserved}}
\qbezier(88,9)(98,9)(108,9)
\put(88,10){\makebox(20,4){\tiny (2,2)}}
\put(108,6){\framebox(12,7){\scriptsize x3}}
\put(104,1){\makebox(20,4){\tiny reversed}}
\qbezier(120,9)(130,9)(140,9)
\put(120,10){\makebox(20,4){\tiny (2,2)}}
\put(140,6){\framebox(12,7){\scriptsize a2}}
\put(136,1){\makebox(20,4){\tiny preserved}}
\end{picture}
\par\smallskip
{\scriptsize\begin{adjustbox}{max width=\linewidth}\begin{tabular}{llllll}
the site & the link & the Gram of the propagation & the signature & the orientation & the boost plane and generator \\ \hline
a1 & a1-x1 (2,2) entered & 4,0,0,0;0,-1,0,0;0,0,-1,0;0,0,0,-1 & 1,3, the ratio -4 & preserved & 4,0;0,-1; 0,1/2;2,0 \\
x1 & a1-x1 (2,2) far & 4,0,0,0;0,-1,0,0;0,0,-1,0;0,0,0,-1 & 1,3, the ratio -4 & reversed & 4,0;0,-1; 0,1/2;2,0 \\
x2 & x1-x2 (2,2) far & 4,0,0,0;0,-1,0,0;0,0,-1,0;0,0,0,-1 & 1,3, the ratio -4 & preserved & 4,0;0,-1; 0,1/2;2,0 \\
x3 & x2-x3 (2,2) far & 4,0,0,0;0,-1,0,0;0,0,-1,0;0,0,0,-1 & 1,3, the ratio -4 & reversed & 4,0;0,-1; 0,1/2;2,0 \\
a2 & x3-a2 (2,2) far & 4,0,0,0;0,-1,0,0;0,0,-1,0;0,0,0,-1 & 1,3, the ratio -4 & preserved & 4,0;0,-1; 0,1/2;2,0 \\
\hline
\multicolumn{6}{l}{the pole of the two-point row at every link 30;30;0;0 over 2304; the order-step plane of the open result 1,0;0,1/4, euclidean at the chain as at the unit} \\
\end{tabular}\end{adjustbox}}
\caption{The signature at propagation along the serial (2,2) chain of the electron at the third step, the sites a1, x1, x2, x3, a2 over the 4 inserted serial (2,2) links: the unit and the three direction steps carried along the chain by the transport identity give the Gram of the propagation 4,0,0,0;0,-1,0,0;0,0,-1,0;0,0,0,-1 at every step, its signature one against three equal the pair-end signature 1,3 at 116 of 116 over the chain states, the ratio -4 at every step, the carried triple reversed in orientation at every odd step (50 of 116); the boost of the chain, the generator 0,1/2;2,0 of the plane 4,0;0,-1 of the propagated Hermitian unit and a carried direction, an isometry of the Gram of the propagation at 116 of 116 (the Lorentzian generator, the light cone of the literature named for the board), while the one-unit rotation of the open result stays Euclidean at the order step at the chain as at the unit. The cone drawn above the chain is the pair of null directions of that plane, the unit entry against the direction entry; the orientation at each site stands beneath it.}\label{fig:III-3-D543402-causal-cone}
\end{figure}

\begin{chbridge}
The Minkowski form has one sign against three, a Lorentz boost is a hyperbolic rotation preserving it, and a Euclidean rotation preserves a definite form instead, so that the signature of the metric fixes which generator propagates \cite{Minkowski1908,Einstein1905}. The derivation supplies the Gram at every step and the two generators at their pairings here, and the signature is named by correspondence.
\end{chbridge}
\begin{chproof}[finite and exact]
The premises are the boost open result at the (2,3) link frame (chapter 3-3, an open result), the (2,2) pair at its Gram and signature (chapter 10-0, a theorem), the Lorentz group at the qubit pair with its brackets table (chapter 3-3, a theorem), the serial insertion row at the chains (chapter 10-6, an open result) and the two-point row at the pole (chapter 10-2, an open result). The chain back to the axiom runs through chapter 10-0 and the framework chapter. By exhaustive enumeration of the one hundred sixteen steps of the thirty-four chain states over the eighty-two serial (2,2) links of the nine recorded webs, the electron, the muon, the tau, the proton, the neutron through the deuteron, the photon word and the universe closure, with the neutron and the tau sink beside, the proton and photon chain links at the zeroth step being (2,3) links with no serial (2,2) link and so empty, the property tested at each step being whether the Gram of the propagation reads one against three and whether each of the two generators is an isometry of its pairing. From the carried unit and the three carried directions at the first inserted site of the electron chain, the transposed triple $0,-\tfrac{i}{2};-\tfrac{i}{2},0$ and $0,\tfrac{1}{2};-\tfrac{1}{2},0$ and $-\tfrac{i}{2},0;0,\tfrac{i}{2}$, by the trace pairing on the link space, we obtain the Gram
\begin{equation}
G = \begin{pmatrix} 4 & 0 & 0 & 0 \\ 0 & -1 & 0 & 0 \\ 0 & 0 & -1 & 0 \\ 0 & 0 & 0 & -1 \end{pmatrix} , \qquad \text{the signature } (1,3) , \quad \text{the ratio } -4 ,
\end{equation}
equal to the (2,2)-pair row at both ends, the signature by the pivot signs at the first engine equal to the leading principal minors at the second at every step. From the Hermitian unit and a direction under the plain trace, by the restriction to their plane, we obtain the boost of the chain
\begin{equation}
P_{L} = \begin{pmatrix} 4 & 0 \\ 0 & -1 \end{pmatrix} , \qquad K = \begin{pmatrix} 0 & \tfrac{1}{2} \\ 2 & 0 \end{pmatrix} , \qquad K^{\mathsf{T}} P_{L} + P_{L} K = 0 ,
\end{equation}
the residual zero and the type boost at one hundred sixteen of one hundred sixteen steps, the mixing block nonzero; and from the anti-Hermitian order step with a direction under minus the trace over six, the order-step plane $1,0;0,\tfrac{1}{4}$, definite, with the one-unit rotation $0,-\tfrac{1}{2};2,0$ its isometry at the residual zero and no isometry of the Gram of the propagation restricted at the residual $0,-4;-4,0$, the open result letter Euclidean. The counts by kind: one against three at one hundred sixteen of one hundred sixteen steps, definite at none, other at none; the boost an isometry at one hundred sixteen of one hundred sixteen; the order-step plane definite with the open result rotation its isometry at one hundred sixteen of one hundred sixteen; the orientation of the carried triple preserved at the even steps and reversed at the odd, fifty of one hundred sixteen, the structure constants at the literal carried basis negated at the odd steps, seventy-eight of ninety at the first inserted site with the zeros standing, and equal to the brackets table at ninety of ninety at the re-oriented basis at every step; the pole at every link at the chain $(2,2)$ with the displacement zero and the two-point function $\tfrac{30 + 30\sqrt{2}}{2304}$ equal to the two-point row at two hundred four of two hundred four. So the signature lives at the choice of the unit and not at the chain: the anti-Hermitian order step pairs with a direction at a definite plane under every trace form, so the open result rotation is a Euclidean rotation at the unit and at every step alike, while the Hermitian unit pairs with a direction at an indefinite plane, so the boost of the chain is the Lorentzian generator at every step, and the transport identity carries the signature along every chain unmoved, which is the theorem. The smallest instance on the page is the first step of the electron chain: the Gram above, the boost preserving the plane $4,0;0,-1$ and the rotation preserving $1,0;0,\tfrac{1}{4}$.

The counter-tests struck the Hermitian unit to the anti-Hermitian order step at the chain, which moved the plane to the definite form and the type to the rotation at every step. A second reversed the orientation at every step, which moved the preserved and reversed counts and left the signature unmoved. A third read the terminal (2,3) link of the proton and photon chains as a chain link, which moved the Gram to $6,0,0,0;0,-1,0,0;0,0,-1,0;0,0,0,-1$ at those two steps with the signature unmoved, and a fourth struck the gate, which left every cell absent. Every control read the same at both engines, and the steps were reproduced at a second engine of a different construction, byte-identical at every run.
\chcert{116}
\end{chproof}
\begin{figure}[h]\centering
\begin{adjustbox}{max width=\linewidth}\begin{tabular}{lllll}\toprule
generator & the plane & the square & the pairing preserved & the letter \\ \midrule
the one-unit rotation $0,-\tfrac{1}{2};2,0$ & $1,0;0,\tfrac{1}{4}$, definite & minus the identity & the Euclidean pairing, residual $0$ & Euclidean \\
the hyperbolic boost $0,\tfrac{1}{2};2,0$ & $4,0;0,-1$, indefinite & the identity & the Lorentzian pairing, residual $0$ & Lorentzian \\
the one-unit rotation at the Gram restricted & $4,0;0,-1$ & minus the identity & not preserved, residual $0,-4;-4,0$ & no isometry \\ \bottomrule
\end{tabular}\end{adjustbox}
\caption{The two generators of the boost open result at their planes and pairings at every step of the serial chain, the rotation Euclidean at the unit and the boost Lorentzian at the chain.}
\end{figure}
\par\medskip\noindent{\itshape From the group the thread passes to the one transformation that mixes order and direction, the boost.}\label{narr:close:3-3}\par
\section{Boosts and the one unit}\label{mod:III-4-boosts-and-the-one-unit}
\chdesig{4-4}{3-4}{III-4-boosts-and-the-one-unit}{III-4-boosts-and-the-one-unit}
\label{narr:open:3-4}Theorem~56.1 places the boost at the one unit, an open result whose parting the box names, resting on the order coordinate and the Lorentz group.
\begin{chtheorem}[The order register at the three directions]\label{box:III-4-boosts-and-the-one-unit:1}
The twirl of the definition \emph{known}, run over the order register and the three directions together at the record of every read, carries the order coordinate into no direction coordinate at the pair ket: the forty-eight matchings of the five coordinates that the construction admits, the twenty-four rotations of the cube on the directions joined with the two signs on the order coordinate, each carry a zero block from the order coordinate into the direction coordinates, so that no record of any read is carried to a record whose direction component is drawn from its order, 0 of 322714 records at 0 of 3072 pairs of a web and a matching over the 64 catalogued webs, the first web named with its forty-eight matchings counted and none carrying. The hand matching that exchanges the order coordinate with one direction carries 319332 of those records and erases the order coordinate at every one of them, but it lies outside the admitted closure, 384 maps against 48; the twirl over the directions alone and the twirl with the reflections carry nothing. The first open result of the previous release at this chapter, that no established rule mixes the order register with the three directions so that the relative sign of the two coefficients is fixed by no rule, is closed at this reading: the nine rules carry a zero block between the order register and the directions because every admitted matching does, and the sign is a property of the twirl and of no rule, the deciding table re-derived in Lean at the forty-eight matchings.
\end{chtheorem}
\chstanding{theorem}
\chdeps{Definitions 1.1 to 1.6, Theorem 50.1 and Theorem 54.1}{no theorem beyond it at this release}
\begin{chwall}[The boost at the one unit]\label{box:III-4-boosts-and-the-one-unit:2}
Two open results stand here. The two ticks are exact, their squared ratio two at the standing unit and one at two admitted alternates, so the integer maps preserving the Lorentzian form are the Pell family at the standing unit and the sign maps alone at the alternates: the boost from the tick ratio is unit-dependent, the one unit the missing object. Inside the anti-Hermitian frame algebra the boost preserves the Euclidean pairing and not the Lorentzian form, the absent object supplied outside it by the first theorem above.
Open are the two results above at their steps, the ratio step of the tick ratio with both faces of the ratio printed and the deciding step inside the frame algebra, with the form of the unit among the words the derivation still owes; the mixing step of the order register is closed by the theorem above.
\end{chwall}
\chstanding{an open result}
\chdeps{Definitions 1.1 to 1.6, Theorem 50.1 and Theorem 54.1}{no theorem beyond it at this release}
The boost the previous releases of the compendium left at one premise is supplied outside the anti-Hermitian frame algebra: the Hermitian order unit carries with the three direction steps the Gram matrix of signature one plus three, the hyperbolic generator solved from it squares to the identity, and its one-parameter group, exact at the rational points of the hyperbola in every direction, is an isometry of the signature form with the unit mixed with the direction at a nonzero block.
The standing boost, applied as a frame map to the four coordinates of every read, the level the first power of the web and the three components the direction functional on the arrow of the read, preserves the shell invariant at every read, the level moved by the direction and the direction by the level wherever either is nonzero, so the mass shell at the tabulated reads is the level set of the invariant.
The one unit tying the two ticks is a theorem of the readout trace, one tick per act, the unit of the order register fixed by the trace of the readout at the one link; the form of the unit, the pairing the standing lattice group admits, stands among the words the derivation still owes.
\begin{chbridge}
A boost of velocity $v$ along the first axis is the hyperbolic rotation of the $(t, x_1)$ plane by the rapidity $\eta$ with $\tanh \eta = v$, the invariant interval preserved, and its generator is the Hermitian element $K_1$ of the Lorentz algebra \cite{Einstein1905,Weinberg1995}; the invariant mass shell of a particle is the level set $E^2 - p^2 = m^2$, carried to itself by every boost \cite{Weinberg1995}. The correspondence is by name: the compendium supplies a joined saturation carrying the hydrogen name, a two-vertex term carrying the Rydberg name, cells at the hydrogen sink carrying the Lamb and hyperfine names, self-terms carrying the coupling names and joined electron sinks carrying the periodic-table name at chapter~\ref{mod:IV-8-ward-identity-and-the-anomaly-sums}, chapter~\ref{mod:V-2-electron-anomalous-dipole-factor-through-order-three}, the chapter titled The muon anomalous dipole factor (chapter~\ref{mod:V-3-muon-anomalous-dipole-factor}, among the instruments), chapter~\ref{mod:V-4-hydrogen-the-rydberg-ladder-and-the-fine-structure}, chapter~\ref{mod:V-5-lamb-shift-and-the-casimir-effect}, chapter~\ref{mod:V-8-photon-docking-the-join-face-and-pair-production}, the chapter titled The mesons and the nucleon at the cap (chapter~\ref{mod:VII-9-mesons-and-the-nucleon-at-the-cap}, among the instruments), chapter~\ref{mod:XII-3-chain-of-pairs-and-the-charged-periodic-sink}, chapter~\ref{mod:XII-5-periodic-table-as-the-joined-electron-sinks} and chapter~\ref{mod:XIV-7-what-remains-the-ungraded-page}, a graph Laplacian over a ladder carrying the curvature name, a register tower carrying the metric name, a non-sink horizon web, a depth exponent carrying the cosmological-term name and a passage of the read carrying the geodesic name at chapter~\ref{mod:I-3-records-reads-and-the-born-count-the-gain-lemma}, chapter~\ref{mod:VIII-1-mass-unit-and-the-first-power}, chapter~\ref{mod:IX-4-imbalance-readout-lifetimes-and-branching}, chapter~\ref{mod:X-1-curvature-as-the-graph-laplacian-washed-over-the-ladder}, chapter~\ref{mod:X-2-field-equation-at-the-counting-face}, chapter~\ref{mod:X-3-horizon-at-every-rung-and-the-horizon-web-as-a-seeded-saturation}, chapter~\ref{mod:X-4-depth-the-growth-order-and-the-termination-wave}, chapter~\ref{mod:X-5-register-tower-and-the-metric}, chapter~\ref{mod:X-6-cosmological-term-at-the-depth-exponent}, the chapter titled The horizon as a non-sink: the Hubble rate and the baryon asymmetry (chapter~\ref{mod:X-8-horizon-as-a-non-sink-the-hubble-rate-and-the-baryon-asymmetry}, among the instruments), the chapter titled The gravitational and cosmological constants table (chapter~\ref{mod:X-9-gravitational-and-cosmological-constants-table}, among the instruments), chapter~\ref{mod:X-10-tick-scale-as-a-coupling-scale} and chapter~\ref{mod:XIV-7-what-remains-the-ungraded-page} and a step generator at the order coordinate, a first-power operator, a direction functional and an office at the common support carrying the energy, momentum and mass names at chapter~\ref{mod:I-3-records-reads-and-the-born-count-the-gain-lemma}, chapter~\ref{mod:I-6-read-at-a-register-the-probe-and-the-measurement}, chapter~\ref{mod:II-6-energy-and-the-hamiltonian-the-first-power-operator-and-the-step}, chapter~\ref{mod:III-1-order-coordinate-and-its-theorem}, chapter~\ref{mod:III-6-p-c-t-and-cpt}, chapter~\ref{mod:IV-6-conservation-as-the-kernel-of-the-move-matrix}, chapter~\ref{mod:V-7-running-of-the-coupling}, the chapter titled The Dirac equation at the dock (chapter~\ref{mod:V-9-dirac-equation-at-the-dock}, among the instruments), chapter~\ref{mod:VIII-1-mass-unit-and-the-first-power}, chapter~\ref{mod:VIII-5-unified-mass-readout}, chapter~\ref{mod:X-2-field-equation-at-the-counting-face}, the chapter titled The gravitational and cosmological constants table (chapter~\ref{mod:X-9-gravitational-and-cosmological-constants-table}, among the instruments), chapter~\ref{mod:XI-1-second-law-as-the-boundary-count-under-acts} and chapter~\ref{mod:XIV-7-what-remains-the-ungraded-page}, and the table named as the shell filling of a Coulomb Hamiltonian under the exclusion principle, the theory named as the metric theory of a spacetime with the Einstein equations and the momentum named as the generator of translations and a label of the Lorentz group correspond to these objects by name.
\end{chbridge}
\begin{chproof}[exact part and census part]
The premises of this derivation are the axiom and the six definitions of chapter~\ref{mod:FRAMEWORK}, the table of the chapter titled The Rosetta table in full (chapter~\ref{mod:I-7-rosetta-table-in-full}, among the instruments), the theorem of chapter~\ref{mod:III-1-order-coordinate-and-its-theorem} and the theorem of chapter~\ref{mod:III-3-lorentz-group-and-the-signature}. The chain of premises runs back to the axiom through those chapters, and the longest such chain in this compendium is twelve chapters long.
The Gram matrix of the Hermitian order unit and the direction steps under the plain trace is the Lorentzian pairing multiplied by three halves and the hyperbolic generator solved from it is exponentiated at the rational points of the hyperbola; the boost gate is applied at the eight exact points of the hyperbola to the four coordinates of every read with the shell invariant evaluated ahead of and beyond the map; and the trace of the readout at the one link fixes the unit of the order register; the theorems are conditional on the readings of the construction and on the open results they rest on, among them the boost from the order register and the momentum as the direction functional.
\textit{The exact part.}
From the Hermitian order unit and the three direction steps under the plain trace, by the Gram matrix, we obtain the Lorentzian pairing multiplied by three halves, the signature one plus three. From that Gram matrix, by solving for the generator that mixes the unit with one direction at a nonzero block, we obtain the hyperbolic generator, which squares to the identity.
\begin{equation}
G = \tfrac{3}{2}\,\operatorname{diag}(4, -1, -1, -1), \qquad K^{2} = 1 .
\end{equation}
From the hyperbolic generator, by exponentiation at the rational points of the hyperbola, we obtain a one-parameter group of isometries of the signature form, exact at every rational point, the unit mixed with the direction at a nonzero block. Applied as a frame map to the four coordinates of every read (the level the first power of the web and the three components the direction functional on the arrow of the read), each boost preserves the shell invariant,
\begin{equation}
\exp(\theta K)^{\top} G \exp(\theta K) = G, \qquad I(\ell, \vec{p}\,) = I(\ell', \vec{p}\,') .
\end{equation}
From the integral form of the Lorentzian pairing at the standing tick ratio, by the fresh enumeration of the integer maps preserving it at the bounds ten, one thousand and one hundred thousand, we obtain four maps at every bound with no mixing, the sign maps alone and no first boost, which names the parting of this open result, the boost from the tick ratio is unit-dependent and the one unit is the missing object,
\begin{equation}
\#\{\text{integer maps of } (4, -1)\} = 4 \ \text{at every bound}, \qquad \#\{\text{mixing maps}\} = 0 .
\end{equation}
\textit{The census part.}
By exhaustive enumeration of the 309 tabulated reads, the boost gate is applied at the eight exact points of the hyperbola to the four coordinates of every read with the shell invariant evaluated ahead of and beyond the map, preserved at every read, the level moved by the direction and the direction by the level wherever either is nonzero, so the mass shell at the tabulated reads is the level set of the invariant. The order register is re-met with the arrow entering a site at every one of the 3382 arrow lines, and the full census (32851 instances, 6296 cuts, 662 fixed rows and 23320 joins) is re-met at both engines. The smallest case, the Gram matrix and the generator at the qutrit-qubit link, is worked above.
The counter-tests alter the objects singly and read what moves, every one of the five at its pinned letter at both engines, and the fifteen controls read the same at both engines. The result was reproduced by a second engine built apart from the first, the two runs equal at every key.
The machine-checked package named at the reproducibility appendix certifies the statements of this chapter under the definitions of this compendium and proves no theorem of the literature.
\chcert{309}
\end{chproof}
\begin{figure}[h]\centering
\setlength{\unitlength}{1pt}
\begin{picture}(260,150)
\qbezier(86.0,75.0)(130.0,75.0)(174.0,75.0)
\put(130.0,83.0){\makebox(0,0){\scriptsize $(2,3)$}}
\put(77.0,75.0){\circle{14}}
\put(77.0,75.0){\makebox(0,0){\footnotesize $3$}}
\put(77.0,85.0){\makebox(0,0)[b]{\scriptsize $q$}}
\put(183.0,75.0){\circle{14}}
\put(183.0,75.0){\makebox(0,0){\footnotesize $2$}}
\put(183.0,85.0){\makebox(0,0)[b]{\scriptsize $b$}}
\end{picture}
\caption{The qutrit-qubit link whose frame algebra carries the Hermitian order unit outside its anti-Hermitian part. The boost of this chapter mixes that unit with a direction step at the signature form.}
\end{figure}
\par\medskip\noindent{\itshape Beyond the boost, the next chapter reads the label that the literature calls spin.}\label{narr:close:3-4}\par
\section{Spin and helicity}\label{mod:III-5-spin-and-helicity}
\chdesig{4-5}{3-5}{III-5-spin-and-helicity}{III-5-D542908-spcf-t1}
\label{narr:open:3-5}Theorem~57.1 reads spin from the open-end label, an open result with its parting named, resting on the rotation group of chapter 3-2.
\begin{chtheorem}[Spin at the differing body]\label{box:III-5-spin-and-helicity:1}
At the first body at which the twice-spin read from the open qubit ends alone parted from the name, the docked up quark at the reading that keeps the nonzero labels, a web of three sites and two links at the genre (2,3) with one open ket end and one open bra end of dimension two and no open qutrit or ququart end, the label content re-built at a second route from its open ends alone, the weight route with no pairwise Clebsch-Gordan rule, reads zero and two at four states with the multiplicities one and one, equal to the rule at that body at every one of its three readings, the low at zero against zero, the nonzero at two against two and the maximal at two against two; the arrow spaces of its two links are printed at the dimension six each, the ket end tensor the bra end, with both weight faces as the catalogue carries them, and the body has no ququart end whose frame could part. The construction of the open-end label therefore stands re-met at its first differing body by a second route, and the parting the previous open result named is a parting from the name alone, the body up at the twice-spin one by name, which differs from the second route at every reading and is an input to no cell; the deciding table is re-derived in Lean.
\end{chtheorem}
\chstanding{theorem}
\chdeps{Definitions 1.1 to 1.6 and Theorem 53.1}{Theorem 58.1, Theorem 59.1, Open result 85.1, Theorem 124.1 and Theorem 143.1}
\begin{chwall}[Spin from the open-end label]\label{box:III-5-spin-and-helicity:2}
Two open results stand here. The twice-spin read from the open qubit ends alone reaches the name at six of the nine bodies and one of the five carriers, the up quark first, its construction at that body re-met by the theorem above. The composition of the doublet objects by the symmetric power per passage reads one doublet at every body and the one at no carrier, posed at two of its steps and re-set at the closed read by the first theorem above. The standing spin reading at the eight networks of the up-type quarks and the carriers equals the census doublet content at seven of the eight, the one parting re-set at the attach site by the last theorem above.
Open are the two results beyond the first at their steps, the amplitude step with both readings printed and the one parting at the carriers, the content step at the differing body closed by the theorem above; the identification of the helicity direction with the momentum is a reading named as a premise.
\end{chwall}
\chstanding{an open result}
\chdeps{Definitions 1.1 to 1.6 and Theorem 53.1}{Theorem 58.1, Theorem 59.1, Open result 85.1, Theorem 124.1 and Theorem 143.1}
At the closed reads the content is the triplet at the W, the Z, the gluon and the photon words and none at the Higgs word, while at the dock reads of the nine chain bodies it is one doublet, so spin is the half of every chain body, the one of every vector carrier and the zero of the Higgs; conditional on the posed mixed half of the gain law and the posed join reading.
The read line entering the qutrit-qubit link at its qutrit end leaves the far qubit open at the generic read with exactly one direction of its traceless triple fixed, so the transverse traceless part at the open read has dimension two, helicity by the literature name.
The exchange of the two copies of every seeded sink is the composition of three tabulated moves, its sign at the joined read minus at every dock join and plus at every other join, so the exchange sign equals minus one to the twice-the-spin at every seeded sink class.
The ququart frame is a qubit tensor a qubit at every ququart site of the construction, and under the standing diagonal $\mathfrak{su}(2)$ the spinor, the frame and the arrow spaces decompose exactly, nothing laid.
Under the standing matching the open qubit-end set of every interaction instance changes by matched ket-bra pairs alone, so the signed qubit-tensor count is conserved: angular momentum is a conserved count of the construction, conditional on the one-act clause.
At the two-dock up network joined to its meson apparatus network the port of the up is the one open far end, and the up spin at that site is the spin doublet of the census at the port.
\begin{chbridge}
A massive particle of spin $s$ carries the irreducible representation of $SU(2)$ of dimension $2s + 1$, the doublet at spin one half \cite{Wigner1959}; a massless particle carries a helicity with two transverse states \cite{Wigner1939}; and the exchange of two identical particles carries the sign plus one at integer and minus one at half-integer spin \cite{Pauli1940}. The correspondence is by name: the compendium supplies a finite re-matching group at the known clause, a representation census at the open ends and the sinks, an attach-site basis carrying the generation name and a sink class carrying the doublet name at chapter~\ref{mod:IV-1-gauge-invariance-as-the-known-clause}, chapter~\ref{mod:IV-3-representations-from-the-open-end-census}, chapter~\ref{mod:IV-4-cartan-weights-and-the-charge}, chapter~\ref{mod:IV-6-conservation-as-the-kernel-of-the-move-matrix}, chapter~\ref{mod:IV-7-noether-at-the-sinks-and-the-kernel-identification}, chapter~\ref{mod:IV-9-one-rule-of-a-read-and-the-word-table}, chapter~\ref{mod:VI-3-doublets-handedness-and-parity-violation}, chapter~\ref{mod:VII-7-quark-content-and-flavour}, chapter~\ref{mod:VIII-4-three-generations-as-attach-sites}, chapter~\ref{mod:XIV-7-what-remains-the-ungraded-page}, a re-matching group with a wash integer, a centralizer weight at the kernel, an enclosed record count carrying the flux name and a counting-face equation carrying the field-equation name at chapter~\ref{mod:I-3-records-reads-and-the-born-count-the-gain-lemma}, chapter~\ref{mod:IV-1-gauge-invariance-as-the-known-clause}, chapter~\ref{mod:IV-4-cartan-weights-and-the-charge}, chapter~\ref{mod:IV-6-conservation-as-the-kernel-of-the-move-matrix}, chapter~\ref{mod:IV-7-noether-at-the-sinks-and-the-kernel-identification}, chapter~\ref{mod:V-1-fine-structure-constant-at-the-one-link}, chapter~\ref{mod:VI-2-w-and-the-z-and-their-masses}, chapter~\ref{mod:X-2-field-equation-at-the-counting-face}, chapter~\ref{mod:X-3-horizon-at-every-rung-and-the-horizon-web-as-a-seeded-saturation}, the chapter titled The gravitational and cosmological constants table (chapter~\ref{mod:X-9-gravitational-and-cosmological-constants-table}, among the instruments), the chapter titled The Chern number (chapter~\ref{mod:XII-4-chern-number}, among the instruments) and chapter~\ref{mod:XIV-7-what-remains-the-ungraded-page}, the Rosetta table of this theory nouns, a pair stabilizer, the cap records as a Weyl basis, the Weyl relation at the uncertainty bound, the Bell bound at the pair ket, the gate set at the Weyl criterion, the wash at a sub-web and the Born count as the dressed measure at chapter~\ref{mod:FRAMEWORK}, chapter~\ref{mod:I-1-axiom-and-the-six-definitions}, chapter~\ref{mod:I-3-records-reads-and-the-born-count-the-gain-lemma}, chapter~\ref{mod:I-4-moves-in-one-table}, chapter~\ref{mod:I-6-read-at-a-register-the-probe-and-the-measurement}, the chapter titled The Rosetta table in full (chapter~\ref{mod:I-7-rosetta-table-in-full}, among the instruments), chapter~\ref{mod:II-1-state-space-from-the-act-outputs}, chapter~\ref{mod:II-2-born-rule-and-unitarity}, chapter~\ref{mod:II-3-uncertainty-relation-from-the-weyl-relation}, chapter~\ref{mod:II-9-qm-corollaries-the-fifteen-entries}, the chapter titled Decoherence and the open system (chapter~\ref{mod:II-10-decoherence-and-the-open-system}, among the instruments), chapter~\ref{mod:III-6-p-c-t-and-cpt}, chapter~\ref{mod:III-7-order-theorem-group-moves-and-transfers}, chapter~\ref{mod:IV-2-gauge-group-as-a-corollary-of-one-construction}, chapter~\ref{mod:V-1-fine-structure-constant-at-the-one-link}, chapter~\ref{mod:VI-2-w-and-the-z-and-their-masses}, chapter~\ref{mod:VI-3-doublets-handedness-and-parity-violation}, chapter~\ref{mod:VIII-2-body-alone-office-rule}, chapter~\ref{mod:XI-4-order-parameter-at-the-stop-condition}, chapter~\ref{mod:XI-7-thermodynamic-remainder-and-the-arrow}, chapter~\ref{mod:XIII-1-closure-as-the-computation}, chapter~\ref{mod:XIII-2-moves-as-a-gate-set-and-the-weyl-criterion}, chapter~\ref{mod:XIII-4-distillation-the-read-as-the-maximal-distillation-and-the-distillation-theorem}, chapter~\ref{mod:XIII-5-cost-of-a-read-and-the-adiabatic-class}, chapter~\ref{mod:XIII-6-holonomy-and-the-record-transport}, chapter~\ref{mod:XIV-7-what-remains-the-ungraded-page}, a step generator at the order coordinate, a first-power operator, a direction functional and an office at the common support carrying the energy, momentum and mass names at chapter~\ref{mod:I-3-records-reads-and-the-born-count-the-gain-lemma}, chapter~\ref{mod:I-6-read-at-a-register-the-probe-and-the-measurement}, chapter~\ref{mod:II-6-energy-and-the-hamiltonian-the-first-power-operator-and-the-step}, chapter~\ref{mod:III-1-order-coordinate-and-its-theorem}, chapter~\ref{mod:III-6-p-c-t-and-cpt}, chapter~\ref{mod:IV-6-conservation-as-the-kernel-of-the-move-matrix}, chapter~\ref{mod:V-7-running-of-the-coupling}, the chapter titled The Dirac equation at the dock (chapter~\ref{mod:V-9-dirac-equation-at-the-dock}, among the instruments), chapter~\ref{mod:VIII-1-mass-unit-and-the-first-power}, chapter~\ref{mod:VIII-5-unified-mass-readout}, chapter~\ref{mod:X-2-field-equation-at-the-counting-face}, the chapter titled The gravitational and cosmological constants table (chapter~\ref{mod:X-9-gravitational-and-cosmological-constants-table}, among the instruments), chapter~\ref{mod:XI-1-second-law-as-the-boundary-count-under-acts} and chapter~\ref{mod:XIV-7-what-remains-the-ungraded-page}, a Higgs attachment at the saturated web, a vacuum expectation row as an office register, a dock share and an office ratio carrying the weak-angle name, sink-basis overlaps carrying the mixing names and a W-word office carrying the Fermi name at chapter~\ref{mod:IV-9-one-rule-of-a-read-and-the-word-table}, chapter~\ref{mod:V-12-koide-class}, chapter~\ref{mod:VI-1-hub-step-as-the-weak-interaction}, chapter~\ref{mod:VI-2-w-and-the-z-and-their-masses}, chapter~\ref{mod:VI-3-doublets-handedness-and-parity-violation}, the chapter titled The charged and neutral currents (chapter~\ref{mod:VI-4-charged-and-neutral-currents}, among the instruments), the chapter titled The Fermi constant and the muon decay (chapter~\ref{mod:VI-5-fermi-constant-and-the-muon-decay}, among the instruments), chapter~\ref{mod:VI-6-weak-angle}, chapter~\ref{mod:VI-7-mixing-the-washed-overlap-and-the-ckm-block}, chapter~\ref{mod:VI-10-higgs-the-closed-hub-read-its-rows-its-mass-and-its-exponent}, chapter~\ref{mod:VIII-1-mass-unit-and-the-first-power}, chapter~\ref{mod:VIII-4-three-generations-as-attach-sites}, chapter~\ref{mod:X-1-curvature-as-the-graph-laplacian-washed-over-the-ladder} and chapter~\ref{mod:XIV-7-what-remains-the-ungraded-page} and a dark-genre readability carrying the confinement name, a star office at the common support carrying the nucleon-mass name, a star read at transfer carrying the form-factor name and joined saturations carrying the hadron and nuclear names at chapter~\ref{mod:I-4-moves-in-one-table}, chapter~\ref{mod:III-6-p-c-t-and-cpt}, chapter~\ref{mod:III-7-order-theorem-group-moves-and-transfers}, chapter~\ref{mod:V-1-fine-structure-constant-at-the-one-link}, chapter~\ref{mod:V-7-running-of-the-coupling}, chapter~\ref{mod:V-10-lepton-ladder-as-hub-assemblies}, chapter~\ref{mod:VI-7-mixing-the-washed-overlap-and-the-ckm-block}, chapter~\ref{mod:VII-1-confinement-as-the-dark-genre-readability}, chapter~\ref{mod:VII-3-nucleon-structure-and-the-form-factor}, chapter~\ref{mod:VII-4-proton-as-a-seeded-sink}, chapter~\ref{mod:VII-5-proton-to-electron-ratio-a-candidate}, chapter~\ref{mod:VII-6-hadron-spectrum-and-the-mesons}, chapter~\ref{mod:VII-7-quark-content-and-flavour}, the chapter titled Chiral breaking and the theta angle (chapter~\ref{mod:VII-8-chiral-breaking-and-the-theta-angle}, among the instruments), the chapter titled The mesons and the nucleon at the cap (chapter~\ref{mod:VII-9-mesons-and-the-nucleon-at-the-cap}, among the instruments), chapter~\ref{mod:IX-6-binding-as-the-pairing-records-the-deuteron-and-the-alpha}, chapter~\ref{mod:IX-7-nuclear-rows-at-the-joined-seeded-web}, chapter~\ref{mod:X-7-dark-matter-as-the-dock-free-sinks}, chapter~\ref{mod:XI-4-order-parameter-at-the-stop-condition}, chapter~\ref{mod:XI-5-phase-transition-as-the-sink-class-change} and chapter~\ref{mod:XIV-7-what-remains-the-ungraded-page}, and the multiplet named as a representation of the symmetry, the flavor named as a member of a multiplet of the symmetry with its own mass term, a carrier as the quantum of a connection field with its spin one, the gluon as the quantum of the color connection field with its self-coupling, the object named as a theorem or a standard object of a Hilbert space with its tensor structure, the momentum named as the generator of translations and a label of the Lorentz group, the W as the massive vector quantum of the broken electroweak theory with its mass from the breaking, the Z as the massive vector quantum of the broken electroweak theory with its mass from the breaking, the Higgs as the scalar quantum of the field whose potential breaks the electroweak symmetry and a meson as a quark-antiquark bound state of the confining theory correspond to these objects by name.
\end{chbridge}
\begin{chproof}[exact part and census part]
The premises of this derivation are the axiom and the six definitions of chapter~\ref{mod:FRAMEWORK}, the table of the chapter titled The Rosetta table in full (chapter~\ref{mod:I-7-rosetta-table-in-full}, among the instruments) and the theorem of chapter~\ref{mod:III-2-rotation-group}. The chain of premises runs back to the axiom through those chapters, and the longest such chain in this compendium is twelve chapters long.
The doublet content of the closed read is counted at every one of the eighteen objects of the construction; the read line is followed through the qutrit end to the far-end triple; the exchange of two copies is built from three moves and its sign computed at two engines; the ququart pair state is written as a tensor of qubit pair states; the open qubit-end set is read at its change at every instance; and the up network is joined to its apparatus, each theorem conditional on the results it names, the posed readings, the one-act clause and the open results it rests on.
\textit{The exact part.}
From the end lines of a web, by listing its open qubit copies with their ket or bra letter, we obtain the open-end content and the signed count. At the electron body the open ends are three qubit copies, none ket and three bra, so the signed count is minus three and the label content is the pair of labels one and three.
\begin{equation}
(n, k, l) = (3, 0, 3), \qquad S = k - l = -3, \qquad \{\text{labels}\} = \{1, 3\} .
\end{equation}
From the label content, by the three posed readings of the label rule (the lowest label, the lowest nonzero label and the largest label), we obtain the twice-spin of the electron as one, one and three at the three readings, equal to the name at the first two and differing at the third.
\begin{equation}
2s_{\mathrm{low}} = 1, \qquad 2s_{\mathrm{lnz}} = 1, \qquad 2s_{\max} = 3, \qquad 2s_{\mathrm{name}} = 1 .
\end{equation}
\textit{The census part.}
By exhaustive enumeration of the sixty-four webs of the construction, the open qubit ends are listed from the end lines at every web (two hundred and nineteen open qubit copies, thirty-three ket and one hundred and eighty-six bra) and re-met at every established table and at the nine bodies. The label rule is posed at its readings and the twice-spin read from the open qubit ends alone is compared by name at every body and carrier, it reaches the name at six of the nine bodies and at one of the five carriers, the up quark the first that parts, which names the parting of this open result. The doublet content of the closed read is counted at every one of the eighteen objects, the read line followed through the qutrit end to the far-end triple, and the exchange sign of two copies computed at two engines. The smallest case, the electron body, is worked above.
The counter-tests alter the objects singly and read what moves. An altered label rule moves fifty-seven cells among them the readings and their consistency. Altered open ends move one hundred and thirty-nine web cells. The qutrit counter-test moves one hundred and thirty-seven. The text, names and reading counter-tests are excluded at the gate as pinned. The six counter-tests stand at their pinned letters and the controls read the same at both engines. The result was reproduced by a second engine built apart from the first, the two runs equal at every key.
The machine-checked package named at the reproducibility appendix certifies the statements of this chapter under the definitions of this compendium and proves no theorem of the literature.
\chcert{219}
\end{chproof}
\begin{figure}[h]\centering
\setlength{\unitlength}{1pt}
\begin{picture}(260,150)
\qbezier(127.2,119.4)(114.4,80.1)(101.6,40.7)
\put(122.0,77.6){\makebox(0,0){\scriptsize $(3,3)$}}
\qbezier(122.7,122.7)(104.8,109.7)(86.9,96.7)
\put(109.5,103.2){\makebox(0,0){\scriptsize $(2,3)$}}
\qbezier(107.8,32.1)(130.0,32.1)(152.2,32.1)
\put(130.0,40.1){\makebox(0,0){\scriptsize $(2,3)$}}
\qbezier(106.1,37.4)(139.6,61.8)(173.1,86.1)
\put(134.9,68.2){\makebox(0,0){\scriptsize $(2,3)$}}
\put(130.0,128.0){\circle{14}}
\put(130.0,128.0){\makebox(0,0){\footnotesize $3$}}
\put(130.0,138.0){\makebox(0,0)[b]{\scriptsize $B$}}
\put(79.6,91.4){\circle{14}}
\put(79.6,91.4){\makebox(0,0){\footnotesize $2$}}
\put(79.6,101.4){\makebox(0,0)[b]{\scriptsize $w_1$}}
\put(98.8,32.1){\circle{14}}
\put(98.8,32.1){\makebox(0,0){\footnotesize $3$}}
\put(98.8,42.1){\makebox(0,0)[b]{\scriptsize $C$}}
\put(161.2,32.1){\circle{14}}
\put(161.2,32.1){\makebox(0,0){\footnotesize $2$}}
\put(161.2,42.1){\makebox(0,0)[b]{\scriptsize $w_2$}}
\put(180.4,91.4){\circle{14}}
\put(180.4,91.4){\makebox(0,0){\footnotesize $2$}}
\put(180.4,101.4){\makebox(0,0)[b]{\scriptsize $a_1$}}
\end{picture}
\caption{The electron body with its open qubit ends, the two qutrit core sites, the two dock qubits and an open qubit leg, each site labeled by its dimension. The open qubit copies carry the labels from which the twice-spin is read.}
\end{figure}
\section{The spin-c group at the qubit pair}\label{mod:III-5-D542908-spcf-t1}
\begin{chtheorem}[The spin-c group at the qubit pair]\label{box:III-5-D542908-spcf-t1:1}
The (2,2) identification group is, at every cell of the comparison, the group the literature names the spin-c group of four dimensions: the seven adopted generators (the three anti-Hermitian letters at each end and the traceless unit, anti-Hermitian at seven of seven, of rank seven) close under the bracket at forty-nine of forty-nine with every bracket printed in the recorded coordinates and the centre one line at the diagonal $1;1;-1;-1$; the group defined by its eight letters of the two blocks under the trace constraint of rank one has a null space of dimension seven, anti-Hermitian and traceless at seven of seven, spanning the recorded space; the Lie map $(X_{1}, X_{2}, t)$ to the block matrix of $X_{1} + t\,i\,1$ and $X_{2} - t\,i\,1$ is injective and bracket-preserving at forty-nine of forty-nine with its image the recorded span and its circle image equal to the traceless unit; the explicit isomorphism is the identity matrix at the recorded order with the brackets equal at forty-nine of forty-nine; the presentation at the exponents $1;1$ has the kernel the roots of unity of order two, and the kernel element $(-I, -I, -1)$ acts on the four-dimensional product as the identity at sixteen of sixteen entries, the unreduced action minus the identity. The pair ket $1,0,0,1$ reshaped is the identity join matrix, the unit of the quaternion algebra closed at sixteen of sixteen; its stabiliser of dimension four is the diagonal su(2) together with the u(1), the u(1) acting on the product as zero at the pair scalar zero; the three rotations of the pair row lie inside the stabiliser at three of three with image zero and the three boosts, Hermitian and outside the identification algebra, move the ket at three of three with the images $0,1,1,0$ and $0,-i,i,0$ and $1,0,0,-1$. The Lorentz algebra at a qubit end, the six lifts with their ninety brackets, is the Lorentzian real form of the complexified pair, its constants equal to the table of the Lorentz-group row at ninety of ninety with the sign minus, while the compact real form parts at eighty-seven of ninety with the first parting at the coefficient three of the bracket of the first two boosts, one against minus one; the one-unit rotation of the boost open result, the generator $0,-\tfrac{1}{2};2,0$ squaring to minus the identity, preserves the compact form and not the Lorentzian, and the hyperbolic generator $0,\tfrac{1}{2};2,0$ squaring to the identity preserves the Lorentzian form, the form of the propagation. The single open qubit end is the $(2,1)$ representation of dimension two with the Clifford pair of the dock open result re-met, the sum $(2,1)+(1,2)$ of dimension four is the register with the four generators satisfying the Clifford relation, and the induced mass matrix of the dock open result, the identity of trace two, carries no induced alpha, so that the parting of that open result is re-met. The first face reads yes at the four comparisons and the seventeen pinned cells, seventeen of seventeen; the second face reads no, the first parting none.
\end{chtheorem}
\chstanding{theorem}
\chdeps{Definitions 1.1 to 1.6, Theorem 35.1, Theorem 54.1, Theorem 57.1 and Open result 57.2 and Theorem 124.1}{no theorem beyond it at this release}
The seven generators are re-derived from the identification algebra, the group is built from its defining blocks, the Lie map, the explicit isomorphism and the presentation are solved and checked bracket by bracket, the stabiliser of the pair ket is read with its rotations and boosts, the Lorentz algebra of the qubit end is placed among the real forms of the complexified pair with the boost open result re-read, and the Dirac form is re-read as the sum of the two ends with the dock open results re-met at their partings.

\begin{chbridge}
The spin-c group of four dimensions is the product of the two special unitary groups of rank one with the circle, divided by the diagonal element minus one, and its spinor representation on four dimensions carries the Dirac matrices, four anticommuting matrices squaring to the identity \cite{Dirac1928,AtiyahBottShapiro1964}. The derivation supplies the isomorphism, the stabiliser, the real form and the Clifford relation here, and the spin-c group is named by correspondence.
\end{chbridge}
\begin{chproof}[finite]
The premises are the stabiliser of the pair ket with its identification algebra (chapter 5-1, an open result), the Lorentz group at the qubit pair with its brackets table (chapter 3-3, a theorem), the boost open result (chapter 3-3, an open result), the dock open result and the dock-register open result at the Clifford pair (chapter 3-5, open results) and the photon row at the centre (chapter 6-7, an open result). The chain back to the axiom runs through chapter 2-1 and the framework chapter. The procedure. The input is the seven recorded generators, the eight letters of the two blocks, the pair ket $1,0,0,1$, the six lifts of the qubit end and the two Clifford letters of the dock open result; the rule solves the isomorphism as a change of basis of rank seven, checks every bracket of the forty-nine pairs at both sides, enumerates the kernel of the presentation by its exponent arithmetic, decomposes the stabiliser of the pair ket into its special unitary and circle parts, places the real forms of the complexified pair against the recorded table, and re-reads the Dirac form at the sum of the two ends. From the eight letters under the trace constraint of rank one, by the null space, we obtain the seven adopted generators,
\begin{equation}
\mathfrak{g} = \mathfrak{su}(2) \oplus \mathfrak{su}(2) \oplus \mathfrak{u}(1) , \qquad \dim \mathfrak{g} = 3 + 3 + 1 = 7 ,
\end{equation}
anti-Hermitian at seven of seven with the trace zero, and from the Lie map by the block matrices, the image the recorded span with the circle image $i,0,0,0;0,i,0,0;0,0,-i,0;0,0,0,-i$ equal to the traceless unit; the explicit isomorphism is the identity at the recorded order, the brackets equal at forty-nine of forty-nine. From the presentation at the exponents one and one, by the exponent arithmetic, the kernel is the roots of unity of order two, and the kernel element $(-I, -I, -1)$ has the block form the identity and the tensor action the identity at sixteen of sixteen entries,
\begin{equation}
(-I) \otimes (-I) \cdot (-1)^{0} = I_{4} ,
\end{equation}
the unreduced action minus the identity, so the kernel acts as the identity on the product and the quotient is the spin-c group. From the pair ket reshaped as the identity join matrix, by the action rule $M \mapsto X_{1} M + M X_{2}^{\mathsf{T}}$ checked at twenty-eight of twenty-eight, the stabiliser has dimension four at the coordinates the rows carry, the diagonal su(2) of dimension three with the u(1) inside, their sum the stabiliser, the u(1) tensor action zero at the pair scalar zero; the quaternion algebra closes at sixteen of sixteen with the pair ket its unit, the rotations lie inside at three of three with image zero and the boosts outside at three of three with the images the rows print. From the six lifts at the qubit end, by the fifteen brackets, the constants equal the brackets table at ninety of ninety with the sign minus, the complexified pair at the real rank twelve; the compact real form of rank six lies inside the identification algebra with its constants against the table at eighty-seven of ninety, the first parting the coefficient three of the first two boosts, one against minus one; the Lorentzian real form of rank six carries the boosts Hermitian and outside the identification algebra with the constants equal at ninety of ninety, its boosts $i$ multiplied by the compact boosts at three of three. By exhaustive enumeration of the two generators of the boost open result at the two pairings, the one-unit rotation $0,-\tfrac{1}{2};2,0$ with the square minus the identity preserves the Euclidean pairing $4,0;0,1$ at the residual zero and not the Lorentzian pairing $4,0;0,-1$ at the residual $0,-4;-4,0$, and the hyperbolic generator $0,\tfrac{1}{2};2,0$ with the square the identity preserves the Lorentzian pairing, the form of the propagation of the pair row; the face predicted for the signature row is named here. From the single open qubit end as the $(2,1)$ representation of dimension two, with the Clifford letters $X = 0,1;1,0$ and $Z = 1,0;0,-1$ of the dock open result re-met, the sum $(2,1)+(1,2)$ of dimension four is the register with $\beta = 1,0,0,0;0,1,0,0;0,0,-1,0;0,0,0,-1$ traceless and squaring to the identity and the three alphas squaring to the identity and anticommuting with beta and among themselves at three of three, the Clifford relation at four generators met; the induced mass matrix of the dock open result, $1,0;0,1$ of trace two, lifts to the sum at trace four and not traceless with the induced alpha absent, the Clifford relation at the induced set unmet, which re-meets the parting of the dock open result; the dock-register open result is re-met at the commutants nine and one. The smallest instance on the page is the kernel element at the product: minus the identity at each factor, the product the identity.

The counter-tests struck the circle factor, which moved the centre and the group off the recorded span. A second put the compact real form in place of the Lorentzian at the propagation, which moved the constants to the parting at the first two boosts. A third read the induced set with an alpha inserted, which met the Clifford relation at the induced set and moved the dock open result off its parting, and a fourth struck the gate, which left every cell absent. Every control read the same at both engines, and the seventeen cells were reproduced at a second engine of a different construction, byte-identical at every run.
\chcert{676}
\end{chproof}
\par\medskip\noindent{\itshape The next chapter turns from the continuous symmetries to the discrete ones.}\label{narr:close:3-5}\par
\section{P, C, T and CPT}\label{mod:III-6-p-c-t-and-cpt}
\chdesig{4-6}{3-6}{III-6-p-c-t-and-cpt}{III-6-D542993-nvli-t1}
\label{narr:open:3-6}Theorem~59.1 reads parity as the improper element, and Open result~60.1 states, as an open result, the discrete laws as the Weyl reflections at the cap, a result that reaches forward to the gate set of Part XV.
\begin{chconditional}[Parity as the improper element]\label{box:III-6-p-c-t-and-cpt:1}
This result is conditional on the axiom text and on the results of the chapters named as its premises in the derivation below, each at its own standing.
The lattice group of the three directions has forty-eight signed permutation matrices, half improper, while the identification group at the qubit end acts through eight proper ones and no identification unitary induces an improper element, so parity is not a symmetry of the records.
\end{chconditional}
\chstanding{a conditional theorem}
\chdeps{Definitions 1.1 to 1.6, Theorem 53.1 and Theorem 57.1 and Open result 57.2}{Open result 60.1, Theorem 61.1, Open result 128.1 and Theorem 258.1}
The CPT element of the theory composes the conjugate basis at the far end, the passing along the pair and the reversal of the read order, and the T element is the CPT element with its improper factor struck, conditional on the carrier-phase order as a datum.
At real weights the arrow transfer operator is even under the reversal, so the T-odd and CP-odd content is zero at every cycle, while at the carrier reading the fourth root per read yields a CP measure, odd under the reversal and even under CPT, conditional on the carrier datum.
The fourth root per read may sit at the near end, the far end or both ends with conjugate roots, and CPT is exact at the conjugate pair alone.
At a qubit end the ordered product of the two reads of a sign record has order four exactly, so the carrier phase has order four and the relative component sign reads minus at every reading with a qubit end.
At the five placements the definitions admit and every exponent of the root, every first-order cell of the two-point result stands unmoved while the second-order orientation cells move: the placement is a convention entering no first-order value.
The polarity of the second-order read is derived at every cut from the read line and the standing slot rule, the CP datum a datum under it; conditional on the slot rule.
A read line entering a dock link at its qutrit core leaves one qubit end open at the generic read, while one entering at the qubit opens none, so handedness is the orientation at the twirl.
Four open results stand here. The read operator couples one orientation of its link alone, the orientation coupled no one rule across the reads. The fourth-root carrier phases are forced by no established cell. The relative sign of the two component arrows is defined at ten of the fifty-four tether links alone, the C-odd set of the cochain empty. The conjugate placement is unique, but every root of unity is exact there, so CPT exactness fixes no value. A further open result placed here, the discrete laws at the reversals, is cited at its own standing.
\begin{chbridge}
Parity is the inversion $x \mapsto -x$, violated by the weak interaction \cite{Lee1956}; charge conjugation carries a particle to its antiparticle, $T$ is antiunitary, and $CPT$ is an exact symmetry of every local quantum field theory \cite{Luders1957,Jost1957,StreaterWightman1964}; $CP$ violation enters through one phase of the mixing matrix \cite{Kobayashi1973}. The correspondence is by name: the compendium supplies a re-matching group with a wash integer, a centralizer weight at the kernel, an enclosed record count carrying the flux name and a counting-face equation carrying the field-equation name at chapter~\ref{mod:I-3-records-reads-and-the-born-count-the-gain-lemma}, chapter~\ref{mod:IV-1-gauge-invariance-as-the-known-clause}, chapter~\ref{mod:IV-4-cartan-weights-and-the-charge}, chapter~\ref{mod:IV-6-conservation-as-the-kernel-of-the-move-matrix}, chapter~\ref{mod:IV-7-noether-at-the-sinks-and-the-kernel-identification}, chapter~\ref{mod:V-1-fine-structure-constant-at-the-one-link}, chapter~\ref{mod:VI-2-w-and-the-z-and-their-masses}, chapter~\ref{mod:X-2-field-equation-at-the-counting-face}, chapter~\ref{mod:X-3-horizon-at-every-rung-and-the-horizon-web-as-a-seeded-saturation}, the chapter titled The gravitational and cosmological constants table (chapter~\ref{mod:X-9-gravitational-and-cosmological-constants-table}, among the instruments), the chapter titled The Chern number (chapter~\ref{mod:XII-4-chern-number}, among the instruments) and chapter~\ref{mod:XIV-7-what-remains-the-ungraded-page}, and a carrier as the quantum of a connection field with its spin one corresponds to these objects by name.
\end{chbridge}
\begin{chproof}[exact part and census part]
The premises of this derivation are the axiom and the six definitions of chapter~\ref{mod:FRAMEWORK}, the table of the chapter titled The Rosetta table in full (chapter~\ref{mod:I-7-rosetta-table-in-full}, among the instruments), the theorem of chapter~\ref{mod:III-2-rotation-group} and the open result of chapter~\ref{mod:III-5-spin-and-helicity}. The chain of premises runs back to the axiom through those chapters, and the longest such chain in this compendium is twelve chapters long.
The lattice group and its induced image are enumerated with the determinant of every identification unitary; the CPT element is composed from its three operations; the arrow operator is read under the reversal and the CP measure computed cycle by cycle at three placements; the Weyl product is computed at every qubit end; the five placements are enumerated at every exponent; the polarity is derived from the read line and the slot rule; and the read line is entered at both ends of a dock link, each theorem conditional on the carrier datum, the slot rule and the open results it rests on.
\textit{The exact part.}
From the three derived directions at the qubit end, by the signed permutations of the three axes preserving the Killing form, we obtain the lattice group of forty-eight elements, the six permutations at the eight sign tuples, half of them proper and half improper by the determinant.
\begin{equation}
|L| = 6 \cdot 8 = 48, \qquad \#\{\det = +1\} = 24, \qquad \#\{\det = -1\} = 24 .
\end{equation}
From the identification group at the qubit end, thirty-two monomial unitaries with fourth-root phases, by conjugation on the three directions, we obtain its induced image, eight signed permutation matrices, every one of determinant one, the kernel the four scalar phases, the identity carrying four preimages among the thirty-two and every improper element of the lattice group carrying none.
\begin{equation} \adjustbox{max width=\dimexpr\linewidth-3em\relax}{$\displaystyle \begin{gathered} |G| = 32, \qquad |\operatorname{im} G| = \tfrac{32}{4} = 8, \qquad \det g = +1 \ \text{for all } g \in \operatorname{im} G, \\ \#\{\text{preimages of an improper element}\} = 0 . \end{gathered}$} \end{equation}
\textit{The census part.}
By exhaustive enumeration of the forty-eight signed permutation matrices of the lattice group of the three directions, each with its determinant, trace, order and membership in the induced image, and of the thirty-two identification unitaries with the determinant of the signed permutation each induces, the improper element is absent from the symmetry group of the records at the directions, the induced image is the eight proper elements and no identification unitary induces an improper one, so parity is not a symmetry of the records. The CPT element is composed from its three operations, the arrow operator read under the reversal and the CP measure computed cycle by cycle at three placements, the Weyl product computed at every qubit end, the five placements enumerated at every exponent, and the read line entered at both ends of a dock link. The smallest case, the identity of the image with its four preimages, is worked above.
The counter-tests alter the objects singly and read what moves. Adjoining the element minus the identity, of determinant minus one with no preimage among the thirty-two identification unitaries, closes a group of order sixteen with eight improper elements inside and moves twenty-six cells. The five counter-tests stand at their pinned letters and the twelve controls read the same at both engines, among them the point-group cells re-met by key. The result was reproduced by a second engine built apart from the first, the two runs equal at every key.
The machine-checked package named at the reproducibility appendix certifies the statements of this chapter under the definitions of this compendium and proves no theorem of the literature.
\chcert{48}
\end{chproof}
\begin{figure}[h]\centering
\setlength{\unitlength}{1pt}
\begin{picture}(260,150)
\qbezier(86.0,75.0)(130.0,75.0)(174.0,75.0)
\put(130.0,83.0){\makebox(0,0){\scriptsize $(2,3)$}}
\put(77.0,75.0){\circle{14}}
\put(77.0,75.0){\makebox(0,0){\footnotesize $3$}}
\put(77.0,85.0){\makebox(0,0)[b]{\scriptsize $q$}}
\put(183.0,75.0){\circle{14}}
\put(183.0,75.0){\makebox(0,0){\footnotesize $2$}}
\put(183.0,85.0){\makebox(0,0)[b]{\scriptsize $b$}}
\end{picture}
\caption{The read link at its qubit end, the identification group of the qubit end acts on the three derived directions of the link, and its induced image contains no improper element.}
\end{figure}
\section{The discrete symmetries as the Weyl reflections at the cap}\label{mod:III-6-D538473-w-mwey}
\begin{chwall}[The discrete laws as the Weyl reflections at the cap]\label{box:III-6-D538473-w-mwey:1}
At every one of the seventeen tabulated capped sinks the cap algebra of every link is the Weyl family of the act on the pair, and its sixteen reflections, the sign patterns at the four letters, are typed at exact matrices over the ring of the twenty-fourth root, all sixteen linear at the (2,2) type, eight linear and eight antilinear at the (2,3) and (2,4) types, and four linear, four antilinear and eight absent at the three remaining types; every reversal of the theory, the five maps of the discrete-laws result with the twelve beside, the hub-leg reversal and the component-sign reversal, is an arrow reversal at its support equal to the axis reflection at every support link, antilinear at a qutrit or ququart end and the identity at a qubit end; which reflections are standing moves is read at the end register, every reflection a move at the (2,2) type and the point reflection alone at the (2,3) and (3,3) types; and against the pattern of the discrete-laws result the sinks meet its entries at six of the fifteen lines, the first parting line the parity map at the core-dock role, read broken here against preserved at the reference, while at the point reflection and the match reading every line meets the record.
Open is the result at its parting, the discrete laws exact where their reflections are standing moves and the first parting line named at both cells; the reversals up to re-matching stand at chapter~\ref{mod:V-1-D538789-w-dlrm}, where every reversal equals a reflection composed with a re-matching move, and the mixing phase at the quark block stands at chapter~\ref{mod:VI-7-mixing-the-washed-overlap-and-the-ckm-block}.
\end{chwall}
\chstanding{an open result}
\chdeps{Definitions 1.1 to 1.6, Theorem 59.1 and Open result 231.1}{no theorem beyond it at this release}
\begin{chbridge}
The discrete symmetries of the Standard Model are parity, charge conjugation and the reversal of the arrow of motion; parity and charge conjugation are each violated by the weak interaction, their product is violated at the mixing phases of the quarks, and the product of all three is an exact symmetry of every local Lorentz-invariant field theory by the CPT theorem, with the reversal represented antiunitarily \cite{Lee1956,Wu1957,Luders1954}. The correspondence is by name: the compendium supplies a web of this theory read at its sinks, a doublet swap sign at the odd cycle, a Dirac form at an established generator and a finite record sum carrying the path-integral name at chapter~\ref{mod:II-1-state-space-from-the-act-outputs}, chapter~\ref{mod:II-5-indistinguishability-and-the-spin-statistics-sign}, chapter~\ref{mod:II-7-path-integral-as-the-finite-record-sum}, chapter~\ref{mod:III-5-spin-and-helicity}, the chapter titled The Dirac equation at the dock (chapter~\ref{mod:V-9-dirac-equation-at-the-dock}, among the instruments), chapter~\ref{mod:VI-2-w-and-the-z-and-their-masses}, chapter~\ref{mod:VI-2b-neutral-lepton-the-hub-unit-with-its-legs-open}, chapter~\ref{mod:VII-5-proton-to-electron-ratio-a-candidate}, chapter~\ref{mod:VIII-1-mass-unit-and-the-first-power}, chapter~\ref{mod:XIV-2-walls-census-in-physics-terms} and chapter~\ref{mod:XIV-7-what-remains-the-ungraded-page} and a Higgs attachment at the saturated web, a vacuum expectation row as an office register, a dock share and an office ratio carrying the weak-angle name, sink-basis overlaps carrying the mixing names and a W-word office carrying the Fermi name at chapter~\ref{mod:IV-9-one-rule-of-a-read-and-the-word-table}, chapter~\ref{mod:V-12-koide-class}, chapter~\ref{mod:VI-1-hub-step-as-the-weak-interaction}, chapter~\ref{mod:VI-2-w-and-the-z-and-their-masses}, chapter~\ref{mod:VI-3-doublets-handedness-and-parity-violation}, the chapter titled The charged and neutral currents (chapter~\ref{mod:VI-4-charged-and-neutral-currents}, among the instruments), the chapter titled The Fermi constant and the muon decay (chapter~\ref{mod:VI-5-fermi-constant-and-the-muon-decay}, among the instruments), chapter~\ref{mod:VI-6-weak-angle}, chapter~\ref{mod:VI-7-mixing-the-washed-overlap-and-the-ckm-block}, chapter~\ref{mod:VI-10-higgs-the-closed-hub-read-its-rows-its-mass-and-its-exponent}, chapter~\ref{mod:VIII-1-mass-unit-and-the-first-power}, chapter~\ref{mod:VIII-4-three-generations-as-attach-sites}, chapter~\ref{mod:X-1-curvature-as-the-graph-laplacian-washed-over-the-ladder} and chapter~\ref{mod:XIV-7-what-remains-the-ungraded-page}, and the quark as a colored quantum of a field confined in a bound state and the matrix named as the unitary rotation between the mass and the weak bases of the fields correspond to these objects by name.
\end{chbridge}
\begin{chproof}[exact part and finite part]
The premises of this derivation are the axiom and the six definitions of chapter~\ref{mod:FRAMEWORK}, the theorem of chapter~\ref{mod:III-6-p-c-t-and-cpt} and the open result of chapter~\ref{mod:XIII-2-moves-as-a-gate-set-and-the-weyl-criterion}. The chain of premises runs back to the axiom through those chapters, and the longest such chain in this compendium is twelve chapters long.
At every capped sink the cap algebra of every link is built as the Weyl family of the act and its sixteen reflections are typed at exact matrices; every reversal of the theory is read as a reflection composed with a move at every support link; the membership of each reflection in the group generated by the one-register gates is decided at the end register; and the pattern of the sinks is compared line by line with the discrete-laws result at its readings, the first parting line named at both cells; the result is an open result, a composition resting on the weakest of its parents, the open results among them taken as open.
\textit{The exact part.}
From the four letters of the cap algebra at the two-qubit type, by the sign patterns at the four letters, we obtain the sixteen reflections as exact matrices over the ring of the twenty-fourth root of unity, every one linear and verified on the pair, with none antilinear and none absent, and from the first reversal map at its support, by the reflection it equals at every link, we obtain the axis reflection along the third letter.
\begin{equation}
|\{\text{reflections at } (2,2)\}| = 2^{4} = 16: \ 16 \ \text{linear}, \ 0 \ \text{antilinear}, \ 0 \ \text{absent}; \qquad R_{\mathrm{map}\,1} = \rho_{Z} .
\end{equation}
\textit{The finite part.}
The procedure. For each of the seventeen tabulated capped sinks and each of its links, build the cap algebra of the link as the Weyl family of the act on the pair, the four letters at the two-qubit type. Form its sixteen reflections, the sign patterns at the four letters, as exact matrices over the ring of the twenty-fourth root of unity, and type each as linear, antilinear or absent by testing it on the pair. Read every reversal of the theory (the five maps of the discrete-laws result with the twelve beside, the hub-leg reversal and the component-sign reversal) as an arrow reversal at its support and find the axis reflection it equals at every support link. Decide the membership of each reflection in the group generated by the one-register gates at the end register, and compare the pattern of the sinks line by line with the discrete-laws result at its readings, naming the first parting line at both cells. The smallest instance is the two-qubit type, the four letters give sixteen reflections, every one linear and verified on the pair, and each of the three axis reflections is a move at the end, the pair and the matching readings, projective and exact alike. At the first reversal map the reflection it equals is the axis reflection along the third letter at every link of its support.
At the types $(2,3)$ and $(2,4)$ eight reflections are linear and eight antilinear, and at the three remaining types four are linear, four antilinear and eight absent. The first parting line of the pattern against the discrete-laws result is named at both cells of the deciding table.
The counter-tests alter the objects singly and read what moves. The first moves five cells, the first parting row and cells, the pattern readings and the typed sinks. Every counter-test stands at its pinned letter at both engines and the fourteen controls read the same at both engines, the tables and entries consumed at their hashes with none parting. The result was reproduced by a second engine built apart from the first, the two runs equal at every key.
\chcert{17}
\end{chproof}
\begin{figure}[h]\centering
\setlength{\unitlength}{1pt}
\begin{picture}(260,150)
\qbezier(86.0,75.0)(130.0,75.0)(174.0,75.0)
\put(130.0,83.0){\makebox(0,0){\scriptsize $(2,2)$}}
\put(77.0,75.0){\circle{14}}
\put(77.0,75.0){\makebox(0,0){\footnotesize $2$}}
\put(77.0,85.0){\makebox(0,0)[b]{\scriptsize $a$}}
\put(183.0,75.0){\circle{14}}
\put(183.0,75.0){\makebox(0,0){\footnotesize $2$}}
\put(183.0,85.0){\makebox(0,0)[b]{\scriptsize $b$}}
\end{picture}
\caption{The two-qubit link at the cap, the smallest instance of this chapter, the cap algebra of the link is the Weyl family of the act on the pair, its sixteen reflections the sign patterns at the four letters.}
\end{figure}
\section{Three involutions and the identity product}\label{mod:III-6-D542993-nvli-t1}
\begin{chtheorem}[Three involutions and the identity product]\label{box:III-6-D542993-nvli-t1:1}
The end swap, the chain reversal and the read reversal are each an involution at the pair ket of every genre and at every link of every recorded web, and their product is the identity on every count: the end swap, the end order reversed and the ket end exchanged with the bra end, acts at the pair ket as the transpose of the coefficient matrix on the support, the far-end operator solved from the equal action equal to the transpose at every unit of the support, four of four, nine of nine and sixteen of sixteen, and the pair-ket expectation equal to the trace at every pair of units, sixteen of sixteen, eighty-one of eighty-one and two hundred fifty-six of two hundred fifty-six; the chain reversal reverses the arrow and reads the chain from its far end; the read reversal exchanges the slot and the far end, carries the exponent to the support less the exponent and flips the sign; each is an involution at six of six genres and, applied twice at the orientation datum of every link of the sixty-four recorded webs, one thousand three hundred thirty-four sites and one thousand six hundred ninety-one links, returns the datum at one thousand six hundred ninety-one of one thousand six hundred ninety-one. The counts recomputed at the image, the register, the office under the body-alone rule and the central scalar as a multiset per web, are unmoved by each of the three at sixty-four of sixty-four webs, while the oriented readings beside move, the end orders and ket ends under the end swap, the arrows under the chain reversal, the slots, exponents and sign sum under the read reversal. The product of the three moves every component of the datum at one thousand six hundred ninety-one of one thousand six hundred ninety-one links, its square is the identity, and it is the identity on every count at sixty-four of sixty-four webs and six of six genres, the product in the other order equal at sixty-four of sixty-four since the three maps act at disjoint components. The three discrete-laws open results re-read at their partings name the junction reversal with the odd-cycle holonomy, a map among the three at none of three, so their partings stand outside this item. The first face reads yes and the second face, the first failing involution or moved count named, reads no, the first parting none.
\end{chtheorem}
\chstanding{theorem}
\chdeps{Definitions 1.1 to 1.6, Theorems 9.1 to 9.7 and Open result 9.8, Theorem 59.1, Theorem 124.1 and Theorem 195.1}{no theorem beyond it at this release}
The three maps are defined at the orientation datum of a link, the end order, the ket end, the arrow, the slot, the exponent and the sign; each is applied at the pair ket of the six genres and at every link of the sixty-four recorded webs, applied twice against the datum, and the counts are recomputed at the image by one function of the undirected structure and compared with the counts at the datum; the product is taken in both orders and the open results of the item are re-read at their partings.

\begin{chbridge}
Charge conjugation, parity and the reversal of motion are each an involution of the state space of a field theory, and their product is a symmetry of every local Lorentz-invariant theory, the theorem of Lueders and Pauli \cite{Luders1954,Pauli1955}. The derivation supplies the three involutions and their identity product on the counts here, and the theorem is named by correspondence.
\end{chbridge}
\begin{chproof}[census]
The premises are the pair ket at the six genres with the transpose row (chapter 1-2, a theorem), the discrete laws at their reversals (chapter 3-6, open results), the body-alone office rule (chapter 8-2, a theorem), the serial insertion row at the register (chapter 10-0, an open result) and the stabiliser and photon rows at the central scalar (chapter 5-1 and chapter 6-7, open results). The chain back to the axiom runs through chapter 1-2 and the framework chapter. By exhaustive enumeration of the one thousand six hundred ninety-one links of the sixty-four recorded webs, every link carrying its orientation datum, the property tested at each link being whether each map applied twice returns the datum and whether every count is unmoved. From the end swap at the pair ket, by the equal action on the ket, we obtain the far-end operator as the transpose of the coefficient matrix on the support,
\begin{equation}
\langle \Phi | A \otimes 1 | \Phi \rangle = \langle \Phi | 1 \otimes A^{\mathsf{T}} | \Phi \rangle = \operatorname{tr} A \, / \, d ,
\end{equation}
equal at every unit of the support at the three dimensions and at every pair of units, an involution at six of six genres; the chain reversal is an involution at six of six; the read reversal carries the exponent to the support less the exponent and flips the sign, its exponent map re-met at the reversal row with the element of the record space equal at the support two and differing at three and four, an involution at six of six. The counts by kind: each map applied twice equal to the datum at one thousand six hundred ninety-one of one thousand six hundred ninety-one; the register, the whole and the cut per link, unmoved at sixty-four of sixty-four webs by each of the three; the office under the body-alone rule at the measure-struck reading unmoved at sixty-four of sixty-four; the central scalar, one at the mixed genres and zero at the homogeneous, as a multiset per web unmoved at sixty-four of sixty-four. From the three maps at disjoint components of the datum, by composition, the product moves every component, the datum image differing at one thousand six hundred ninety-one of one thousand six hundred ninety-one, its square the identity and the product the identity on every count at sixty-four of sixty-four webs and six of six genres, the other order equal at sixty-four of sixty-four. The three open results re-read: the dock-face reversal at every (2,3) link an instance of the end swap at the docks and the arrow reversal an instance of the chain reversal, preserved at six of six cells, and the first-parting rows of the re-matching, Weyl-reflections and reversals open results naming the junction reversal at three of three, a map among the three involutions at none of three, so that the partings of the three open results are re-named outside this item. The smallest instance on the page is the end swap at the (2,2) pair ket: the transpose of the identity join matrix the identity, the ket unmoved.

The counter-tests struck the transpose at the end swap, the far-end operator read as the operator itself, which moved the equal action at the mixed genres. A second took the maps in the other order, pinned to the same, which left every count unmoved. A third struck one link from a web, which moved its register, and a fourth struck the gate, which left every cell absent at both engines. Every control read the same at both engines, and the counts were reproduced at a second engine of a different construction, byte-identical at every run.
\chcert{1691}
\end{chproof}
\par\medskip\noindent{\itshape With the reflections read, the next chapter returns to the order coordinate and proves its theorem.}\label{narr:close:3-6}\par
\section{The order theorem: group moves and transfers}\label{mod:III-7-order-theorem-group-moves-and-transfers}
\chdesig{4-7}{3-7}{III-7-order-theorem-group-moves-and-transfers}{III-7-order-theorem-group-moves-and-transfers}
\label{narr:open:3-7}Theorem~62.1 shows that the group moves preserve the order parameter and the transfers raise it, closing the circle the Part opened at chapter 3-1.
\begin{chtheorem}[The group moves preserve the order parameter and the transfers raise it]\label{box:III-7-order-theorem-group-moves-and-transfers:1}
Over every edge of the standing move tables, the gate edges of the step, holonomy, pairing and wash results and the record edges of the closures and the sink edge sets, the order parameter (the order parameter W), the product over the links of a web of the dimension of the common support, is preserved by the group moves, the step and the holonomy, at every gate edge by construction, a gate changing no dimension, and raised by every transfer (the transfer), the attach, the cut, the hub-forming, the hub steps and the joins, at every read record edge with none non-rising, the increment per kind the ratio census; the reverse-marked moves lower it at every read edge, so every kind is partitioned into preserving, raising and lowering by the data alone, with no mixed kind; every count is a count over the enumerated part carrying the caveat of its closure, and the theorem rests on the weakest of its parents, an open result among them.
\end{chtheorem}
\chstanding{a theorem}
\chdeps{Definitions 1.1 to 1.6 and Theorem 50.1}{Theorem 147.1}
Two open results stand here. The depth transfer at the cross join, read as a census of the construction over the components, has its excess at the other link, the step multi-rung and one-way, a statement at the construction rather than at the sinks, conditional on the one-act clause. The depth of a tabulated network is not a function of its register count, the first parting network at depth one with the register count eight against a network at depth zero at that value, the depth metric (a metric on a lattice) of the gravity chapters standing at its established results.
\begin{chbridge}
An order parameter of a phase is a quantity preserved by the symmetry moves of the phase and changed across a transition, and a conserved quantity of a dynamics is one its generators leave fixed \cite{Landau1937,Noether1918}; in a Markov chain on a finite state space a monotone functional, one that never decreases along any transition, orders the reachable states \cite{Kemeny1960}. The correspondence is by name: the compendium supplies a Higgs attachment at the saturated web, a vacuum expectation row as an office register, a dock share and an office ratio carrying the weak-angle name, sink-basis overlaps carrying the mixing names and a W-word office carrying the Fermi name at chapter~\ref{mod:IV-9-one-rule-of-a-read-and-the-word-table}, chapter~\ref{mod:V-12-koide-class}, chapter~\ref{mod:VI-1-hub-step-as-the-weak-interaction}, chapter~\ref{mod:VI-2-w-and-the-z-and-their-masses}, chapter~\ref{mod:VI-3-doublets-handedness-and-parity-violation}, the chapter titled The charged and neutral currents (chapter~\ref{mod:VI-4-charged-and-neutral-currents}, among the instruments), the chapter titled The Fermi constant and the muon decay (chapter~\ref{mod:VI-5-fermi-constant-and-the-muon-decay}, among the instruments), chapter~\ref{mod:VI-6-weak-angle}, chapter~\ref{mod:VI-7-mixing-the-washed-overlap-and-the-ckm-block}, chapter~\ref{mod:VI-10-higgs-the-closed-hub-read-its-rows-its-mass-and-its-exponent}, chapter~\ref{mod:VIII-1-mass-unit-and-the-first-power}, chapter~\ref{mod:VIII-4-three-generations-as-attach-sites}, chapter~\ref{mod:X-1-curvature-as-the-graph-laplacian-washed-over-the-ladder} and chapter~\ref{mod:XIV-7-what-remains-the-ungraded-page} and a graph Laplacian over a ladder carrying the curvature name, a register tower carrying the metric name, a non-sink horizon web, a depth exponent carrying the cosmological-term name and a passage of the read carrying the geodesic name at chapter~\ref{mod:I-3-records-reads-and-the-born-count-the-gain-lemma}, chapter~\ref{mod:VIII-1-mass-unit-and-the-first-power}, chapter~\ref{mod:IX-4-imbalance-readout-lifetimes-and-branching}, chapter~\ref{mod:X-1-curvature-as-the-graph-laplacian-washed-over-the-ladder}, chapter~\ref{mod:X-2-field-equation-at-the-counting-face}, chapter~\ref{mod:X-3-horizon-at-every-rung-and-the-horizon-web-as-a-seeded-saturation}, chapter~\ref{mod:X-4-depth-the-growth-order-and-the-termination-wave}, chapter~\ref{mod:X-5-register-tower-and-the-metric}, chapter~\ref{mod:X-6-cosmological-term-at-the-depth-exponent}, the chapter titled The horizon as a non-sink: the Hubble rate and the baryon asymmetry (chapter~\ref{mod:X-8-horizon-as-a-non-sink-the-hubble-rate-and-the-baryon-asymmetry}, among the instruments), the chapter titled The gravitational and cosmological constants table (chapter~\ref{mod:X-9-gravitational-and-cosmological-constants-table}, among the instruments), chapter~\ref{mod:X-10-tick-scale-as-a-coupling-scale} and chapter~\ref{mod:XIV-7-what-remains-the-ungraded-page}, and the W as the massive vector quantum of the broken electroweak theory with its mass from the breaking, the object named as a structure of a spacetime geometry and the theory named as the metric theory of a spacetime with the Einstein equations correspond to these objects by name.
\end{chbridge}
\begin{chproof}[census]
The premises of this derivation are the axiom and the six definitions of chapter~\ref{mod:FRAMEWORK}, the table of the chapter titled The Rosetta table in full (chapter~\ref{mod:I-7-rosetta-table-in-full}, among the instruments) and the theorem of chapter~\ref{mod:III-1-order-coordinate-and-its-theorem}. The chain of premises runs back to the axiom through those chapters, and the longest such chain in this compendium is twelve chapters long.
Every edge of the move tables is parsed from the parent results, the order parameter is computed at both ends of every edge and re-met at every web of the closures, the difference is signed and the kinds partitioned by the sign alone, at two engines equal at every cell; the theorem is conditional on its parents at their own standing, among them the gate set, the move graph at the sinks and the seeded saturations, and on the glossary.
By exhaustive enumeration of the 759908 edges of the standing move tables (the gate edges of the step, holonomy, pairing and wash results, the record edges of the closures and the sink edge sets, every edge parsed from its parent result with its move kind, its forward or inverse letter and the order parameter at both ends, the order parameter the product over the links of a web of the dimension of the common support), the sign of the difference is read at every edge and the kinds are partitioned by the sign alone, the group moves, the step and the holonomy, preserve the order parameter at every gate edge by construction, a gate changing no dimension, and every transfer (the attach, the cut, the hub-forming, the hub steps and the joins) raises it at every read record edge with none non-rising, the increment per kind the ratio census, while the reverse-marked moves lower it at every read edge. The edges at an unspelled end are refused as unread, the refusal a parting cell and never a halt. The smallest cases are the gate edge of the step at the one-link web of the two-qubit type, whose order parameter two is unchanged, and the first attach edge of the sink edge set at its record face, which carries the order parameter from two to four, the ratio two and the spectrum raised by one (2,3) link:
\begin{equation} W(\text{step}) : 2 \to 2, \ \text{equal}; \qquad W(\text{attach}) : 2^{1}3^{0} \to 2^{2}3^{0}, \ \text{rising}, \ \text{ratio } 2, \ \Delta\sigma = +(2,3) . \end{equation}
The counter-tests alter the objects singly and read what moves. Swapping the two levels, the gate-level kinds against the transfers, moves the deciding cell, the face, the partings and the preservation counts. The three counter-tests stand at their pinned letters and the sixteen controls read the same at both engines, among them the tables and entries consumed at their hashes and the parser guard. The result was reproduced by a second engine built apart from the first, the two runs equal at every cell.
The machine-checked package named at the reproducibility appendix certifies the statements of this chapter under the definitions of this compendium and proves no theorem of the literature.
\chcert{759908}
\end{chproof}
\par\medskip\noindent{\itshape One chapter of the Part remains, and it is counted empty.}\label{narr:close:3-7}\par
\section{The instrument tables of this Part}\label{mod:IV-L-the-instrument-tables-at-the-repository}
\chdesig{4-L}{}{IV-L-the-instrument-tables-at-the-repository}{IV-L-the-instrument-tables-at-the-repository}
\begin{chtable}[The instrument tables of this Part]\label{box:IV-L-the-instrument-tables-at-the-repository:1}
This Part lists 1 instrument chapter, printed among the instruments at the end of this compendium as chapter~\ref{mod:III-8-causal-structure-and-general-covariance}, and the table beneath names each one with its key cells.
\end{chtable}
\chstanding{instrument}
\chdeps{Definitions 1.1 to 1.6}{no theorem beyond it at this release}
The listing moves no cell and states no result of its own. Each instrument stands at the repository as its chapter stood at the previous release of this compendium, and its tables are read there.
\begin{chbridge}
The tables listed here are instruments of this compendium, each listed under the title of its chapter.
\end{chbridge}
\begin{chproof}[instrument]
The listing in full. For each instrument chapter, printed among the instruments at the end of this compendium, the table gives its title, the result boxes it carries, the standing of each box with its citation, and the count of objects its certificate sentence names where it prints one.
{\small
\begin{longtable}{p{38mm}p{63mm}p{30mm}p{14mm}}
\hline
chapter & results & standing & objects \\
\hline
\endhead
Causal structure and general covariance & The chapter counted empty & instrument & -- \\
\hline
\end{longtable}
}
The rule that produced it. Every chapter of this Part whose first standing beyond a placeholder reads instrument is printed among the instruments at the end of this compendium with its tables, and its key cells are the titles of its boxes, its standing lines and its certificate count, read from its text as it stood at the previous release.
No derivation stands in this chapter; the tables stand at the repository.
\emph{Check.} Every file listed was re-hashed at the repository and read the same as its recorded hash, and the listing was re-cut from the files by a second route.
Nothing stands open at this listing. At this release two results of this Part are re-graded at their chapters: the boost at the one unit (the chapter 4-4), its first open result closed as a theorem at the order register and the three directions; spin and helicity (the chapter 4-5), its first open result closed as a theorem at the differing body.
\end{chproof}
\label{chend:IV-L-the-instrument-tables-at-the-repository}
